\documentclass[a4paper]{article}

% ---------- Encoding, lingua, font ----------
\usepackage[utf8]{inputenc}
\usepackage[T1]{fontenc}
\usepackage{textcomp}
\usepackage[italian]{babel}
\usepackage{algorithm}
\usepackage{algpseudocode}
\renewcommand{\familydefault}{\sfdefault}
\emergencystretch=1.5em
\newcommand{\Rep}{\mathrm{Rep}}
\newcommand{\En}{\mathrm{En}}

% ---------- Colori (UNA sola volta, prima di tikz/tcolorbox) ----------
\usepackage[dvipsnames,table]{xcolor}

% ---------- Matematica ----------
\usepackage{amsmath, amssymb, amsfonts, amsthm}
\usepackage{mathtools}
\DeclarePairedDelimiter\ceil{\lceil}{\rceil}
\DeclarePairedDelimiter\floor{\lfloor}{\rfloor}
\usepackage{cancel}
\usepackage{gensymb}

% ---------- Tabelle, figure, utilità ----------
\usepackage{booktabs, xltabular, tabularx, makecell}
\usepackage{subcaption}
\usepackage{float}
\usepackage{graphicx}
\usepackage{import}
\usepackage{pdfpages}
\usepackage{xifthen}
\pdfminorversion=7
\pdfsuppresswarningpagegroup=1
\graphicspath{{./figures/}}

\renewcommand{\arraystretch}{1.5}
\newcolumntype{Y}{>{\raggedright\arraybackslash}X}

% ---------- TikZ / pgfplots ----------
\usepackage{tikz}
\usetikzlibrary{shapes, arrows, automata, petri, decorations.markings,
                decorations.pathreplacing, positioning, calc, quotes}
\usepackage{circuitikz}
\usepackage{forest}
\usepackage{pgfplots}
\pgfplotsset{compat=1.18, width=10cm}

\tikzset{
    place/.style={circle, thick, draw=black, minimum size=6mm},
    state/.style={circle, thick, draw=black, fill=white, minimum size=6mm},
}

% ---------- Palette (solo quella nuova) ----------
\definecolor{thmcol}{HTML}{2563EB}   % blu     - teoremi
\definecolor{lemcol}{HTML}{7C3AED}   % viola   - lemmi
\definecolor{defcol}{HTML}{059669}   % verde   - definizioni
\definecolor{excol}{HTML}{EA580C}    % arancio - esempi
\definecolor{remcol}{HTML}{6B7280}   % grigio  - osservazioni

% ---------- Listings (riusa la palette nuova) ----------
\usepackage{listings}
\lstdefinestyle{mystyle}{
    backgroundcolor=\color{remcol!8},
    commentstyle=\color{defcol},
    keywordstyle=\color{thmcol},
    numberstyle=\tiny\color{remcol},
    stringstyle=\color{lemcol},
    basicstyle=\ttfamily\footnotesize,
    breakatwhitespace=false, breaklines=true, captionpos=b,
    keepspaces=true, numbers=left, numbersep=5pt,
    showspaces=false, showstringspaces=false, showtabs=false, tabsize=2
}
\lstdefinestyle{terminal}{
    backgroundcolor=\color{black},
    basicstyle=\ttfamily\color{green},
    frame=single, breaklines=true, showstringspaces=false
}
\lstset{style=mystyle}

% ---------- Listings (riusa la palette nuova) ----------
\usepackage{listings}
\lstdefinestyle{mystyle}{
    backgroundcolor=\color{remcol!8},
    commentstyle=\color{defcol},
    keywordstyle=\color{thmcol},
    numberstyle=\tiny\color{remcol},
    stringstyle=\color{lemcol},
    basicstyle=\ttfamily\footnotesize,
    breakatwhitespace=false, breaklines=true, captionpos=b,
    keepspaces=true, numbers=left, numbersep=5pt,
    showspaces=false, showstringspaces=false, showtabs=false, tabsize=2
}
\lstdefinestyle{terminal}{
    backgroundcolor=\color{black},
    basicstyle=\ttfamily\color{green},
    frame=single, breaklines=true, showstringspaces=false
}
\lstset{style=mystyle}

\lstdefinestyle{ocamlpseudo}{
  basicstyle=\ttfamily\small,
  columns=fullflexible,
  keepspaces=true,
  showstringspaces=false,
  xleftmargin=0pt
}

% ---------- Teoremi, definizioni, esempi (tcolorbox) ----------
\usepackage[most]{tcolorbox}
\tcbuselibrary{theorems}

\tcbset{
  thmstyle/.style={
    enhanced, breakable, sharp corners,
    boxrule=0pt, frame hidden,
    borderline west={3pt}{0pt}{#1},
    colback=#1!6,
    colbacktitle=#1!14, coltitle=#1!60!black,
    fonttitle=\bfseries,
    titlerule=0pt,
    separator sign none,
    description delimiters parenthesis,
    before skip=12pt, after skip=12pt,
    left=8pt, right=8pt, top=4pt, bottom=4pt,
    overlay unbroken and first={\tikzset{external/export=false}},
    overlay middle and last={\tikzset{external/export=false}},
  }
}

\newtcbtheorem[number within=section]{theorem}{Teorema}{thmstyle=thmcol}{thm}
\newtcbtheorem[use counter from=theorem]{lemma}{Lemma}{thmstyle=lemcol}{lem}
\newtcbtheorem[use counter from=theorem]{definition}{Definizione}{thmstyle=defcol}{def}
\newtcbtheorem[use counter from=theorem]{example}{Esempio}{thmstyle=excol}{ex}
\newtcolorbox{remark}[1][]{thmstyle=remcol, title={Osservazione}, #1}

% ---------- Esercizi (mdframed) ----------
\usepackage{mdframed}
\theoremstyle{definition}
\newmdtheoremenv[
    linecolor=gray, leftmargin=0, rightmargin=0,
    innertopmargin=8pt, innerbottommargin=8pt
]{exercise}{Esercizio}[section]

% ---------- Figure svg ----------
\newcommand{\incfig}[1]{%
    \def\svgwidth{\columnwidth}
    \import{./figures/}{#1.pdf_tex}
}

% ---------- Hyperref (sempre per ultimo, prima di externalize) ----------
\usepackage{hyperref}
\hypersetup{colorlinks=false}

% ---------- Externalize ----------
\usetikzlibrary{external}
\tikzexternalize[prefix=figures/tikz/, optimize=false]
\tikzexternaldisable


\begin{document}
\begin{titlepage}
	\begin{center}
		\vspace*{1cm}

		\Huge
		\textbf{Planning \& Reinforcement Learning}

		\vspace{0.5cm}
		\LARGE
		UniVR - Dipartimento di Informatica

		\vspace{1.5cm}

		\textbf{Tommaso Toffali}

		\vfill


		\vspace{0.8cm}


		1° Semestre 2026/2027
    \end{center}
\end{titlepage}

\tableofcontents
\pagebreak
\part{Classical Planning}
\section{Introduction to Planning}
\subsection{Historical Context and Motivation}
Planning has been a central topic in artificial intelligence since the inception
of the field at the \textbf{Dartmouth Conference in 1956}, where the foundings
fathers of AI first gathered to define the scope and ambitions of this new discipline.
Planning represent one of the foundamental capabilities required for intelligence behavior:
the ability to reason about sequences of actions that achieve desired goals.
\vspace{1em}
\noindent
\newline
The motivation for studying automated planning are manifold:
\begin{itemize}
    \item \textbf{Autonomous agents}: Enabling robots, software agents, and autonomouse system 
        to make decision and act independently in complex environemnts
    \item \textbf{Resource optimization}: Finding optimal or near-optimal stragies for Resource
        allocation, scheduling, and logistics
    \item \textbf{Scientific applications}: Modeling and solving problems in domain such as space 
        exploration, manufacturing, and healtcare
    \item \textbf{Theoretical foundations}: Understanding the computional complexity and formal properties
        of reasoning about action and change
\end{itemize}
\subsection{The planning Cycle}
The complete planning process involves several interconnected pahses that form a continuos cycle:
\begin{enumerate}
    \item \textbf{Plan Generation}: Given a description of the current state, available actions, and desired
            goals, synthesize a sequence of actions (a plan) that transforms the inital state into a goal state.
            This phase involves:
            \begin{itemize}
                \item Analyzing the initial state to understand what is currently true 
                \item Identifying the gap between the current state and the goal state
                \item Searching throungh the space of possible action sequences 
                \item Evaluating candidate plans for correctness and optimality 
                \item Selecting the best plan according to specified criteria (e.g., shortes length, minimum cost, fastest execution)
            \end{itemize}
    \item \textbf{Plan Deployment}: Execute the generated plan in the actual environment, monitoring the execution
            to detect deviations from expected behavior. This phase includes:
            \begin{itemize}
                \item Translating abstract plan actions into concrete executable commands 
                \item Monitoring sensors and feedback to verify that actions have the expected effects
                \item Detecting discrepancies between predicted and actual states
                \item Maintaining a record of executed actions for debugging and learning
            \end{itemize}
    \item \textbf{Replanning}: When execution monitoring detects that the current plan is no longer valid
            (due to unexpected events, actions failures, or environmental changes), generate a new plan
            or repair the existing one. Replanning strategies includes:
            \begin{itemize}
                \item \textbf{Plan repair}: Modify the existing plan minimally to accomodate the new situation 
                \item \textbf{Replan from scratch}: Generate an entirely new plan from the current state
                \item \textbf{Contingency planning}: Use pre-computed alternative plans for anticipated failures
            \end{itemize}
\end{enumerate}
This cycle reflects the reality that planning systems must operate in dynamic, partially unpredictable environments where
initial plans may need adaptation. The cyle continues iteratively until the goal is achived or deemed unreachable.
\subsubsection{Challenges in Real-World Planning}
Real-World deployment of planning systems faces several challengs:
\begin{itemize}
    \item \textbf{Excution uncertainty}: Actions may fail or have unexpected outcomes
    \item \textbf{Incomplete information}: The planner may not have complete knowledge of the environment
    \item \textbf{Dynamic environments}: The world may change while the plan is beign executed
    \item \textbf{Computional constraints}: Planning must often occur in real-time with limited computational resoruces 
    \item \textbf{Plan quality vs. planning time}: Trade-off between finding optimal plans and responding quickly
\end{itemize}
These challenge motivate the development of robust planning algorihms that cna handle unvertainty, adapt to changes,
and operate efficiently under resource constraints.
\section{Formalization of the Planning Problem}
\subsection{State Transition Systems}
The mathematical foundation of planning rests on the concept of state transition system, which provides a formal model of 
how actions transform states.

\begin{definition}{State Transition System}{State Transition System}
A state transition system is a tuple \( \Sigma = (S,A,\gamma) \) where:
\begin{itemize}
    \item \( S \) is a finite or countably infinite set of \textbf{states}
    \item \( A \) is a finite or countably infinite set of \textbf{actions}
    \item \( \gamma : S \times A \rightarrow S \) is a \textbf{state transition function} that maps a state and an action to a resulting state
\end{itemize}
\end{definition}
The transition function \( \gamma(s,a) = s' \) specifies that executing action \( a \) in state \( s \) result in state
\( s' \). This function encodes the \textbf{dynamics} of the domain - how the world changes in response to actions.
\subsubsection{Properties of State Transition System}
State transition systems can be characterized by serveral important properties:
\begin{itemize}
    \item \textbf{Determinism}: In a deterministic system, \( \gamma \) is a function - each state-action pair leads to 
        exactly one successor state. In non-deterministic systems, multiple outcomes are possible.
    \item \textbf{Reachability}: A state \( s' \) is reachable from state \( s \) if there exists a sequence of actions 
        that transform \( s \) into \( s' \). The reachable state space from an initial state is often much smaller thant the totale
        state sapce.
    \item \textbf{Reversibility}: An action is reversible if there exists another action that undoes its effects. Many real-world
        actions are irreversible (e.g., breaking an object)
    \item \textbf{State space structure}: The connectivity and topology of the sate space significanyly affect planning complexity. Highly connected spaces 
        may have many solution paths, while sparse spaces may have few or no solutions.
\end{itemize}
\subsection{Classical Planning: Foundamental Assumptions}
Classical planning makes several simplyfying assumptions that restrict the class of the problems
considered but enable efficient algorithmic solutions.
\begin{definition}{Classical Planning Framework}{Classical Planning Framework}
The classical planning framework relies on six key assumptions:
\begin{enumerate}
    \item Finitely many states
    \item Finitely many actions
    \item Deterministic transition function
    \item Full observability
    \item Instantaneous actions
    \item No exogenous events
\end{enumerate}
\end{definition}
\begin{enumerate}
    \item \textbf{Finitely many states}: The state space \( S \) is finite, allowing exhaustive search techniques.
        This assumption ensures that the planning is decidable and that search alogirhms will terminate
        \vspace{1em}
        \noindent
        \newline
        \textit{Jusitfication}: While real-world domains may have infinite state spaces (e.g. continuos 
        variables), finite approximations are often sufficient for pratical purposes.
    \item \textbf{Finitely many actions}: The action space \( A \) is finite, ensuring decidability. Each state
        has a finite branching factor in the state space graph.
        \vspace{1em}
        \noindent
        \newline
        \textit{Jusification}: Even in complex domains, the number of distinct action type is 
        typically manageable, though the number of ground actions (instantieted with specific objects)
        may be LARGE
    \item \textbf{Deterministic transition function}: For every state \( s \) and action \( a \), there 
        is exactlu one resulting state \( \gamma(s,a) \). Actions have predictable, certain outcomes.
        \vspace{1em}
        \noindent   
        \newline
        \textit{Justification}: Determinism simplified reasoning about action effects. Non-deterministic 
        planning requires more complex formalism (e.g., Markov Decision Processes)
    \item \textbf{Full observability}: The planner has complete knowledge of the current state at all times.
        There is no uncertainty about which state the system occupies.
        \vspace{1em}
        \noindent
        \newline
        \textit{Justification}: Full observability eliminates the need for belief state reasoning and sensing actions.
        Partial observability requires more sophisticated planning approaches.
    \item \textbf{Instantaneous actions}: Actions have no duration, they occur instantaneously. Temporal 
        reasoning is not required.
        \vspace{1em}
        \noindent
        \newline
        \textit{Justification}: Ignoring action durations simplifies the planning model. Temporal planning
        extends classical planning to handle durations and concurrent actions.
    \item \textbf{No exogenous events}: Only the planning agent can change the world through deliberate actions. 
        The environment remains static unless acted upon.
        \vspace{1em}
        \noindent
        \newline
        \textit{Justification}: Exogenous events (e.g., other agents, natural processes) introduce additional complexity.
        Classical planning assumes a static world between actions.
\end{enumerate}
These assumptions, while restrective, capture a significant and important class of planning problems and provide a foundation
for understanding more complex, realistic planning scenarios.
\subsubsection{Relaxing Classical Assumptions}
Modern planning research has devoloped extensions that relax these assumptions:
\begin{itemize}
    \item \textbf{Probabilistic planning}: Handles non-deterministic actions with probability 
        distributions over outcomes
    \item \textbf{Conformant planning}: Plans under partial observability without sensing 
    \item \textbf{Contingent planning}: Generates conditional plans that include sensing actions 
    \item \textbf{Temporal planning}: Reasons about actions durations and temporal constraints
    \item \textbf{Multi-agent planning}: Coordinates plans among multiple agents
\end{itemize}
\subsection{The Planning Problem}
Given a state transition system \( \Sigma = (S,A,\gamma) \) we can now formally define
the planning problem.
\begin{definition}{Planning Problem}{Planning Problem}
A planning problem is a tuple \( P = (\Sigma, s_0, G) \) where:
\begin{itemize}
    \item \( \Sigma = (S,A,\gamma) \) is a state transition system
    \item \( s_0 \in S \) in the \textbf{initial state}
    \item \( G \subseteq S \) is a set of \textbf{goal states}
\end{itemize}
\end{definition}
Alternatively, goals can be specified as a \textbf{goal formula} \( g \), a logical 
expression that is satisfied by exactly those states in \( G \). This allows more compact and expressive goal 
specifications.
\subsubsection{Goal specification}
Goals can be specified in several ways:
\begin{itemize}
    \item \textbf{Explicit goal states}: Enumerate the set \( G \) of accetable final states.
        This is impratical for large state spaces.
    \item \textbf{Goal formula}: A logical formula \( g \) such that 
        \( G = {s \in S \;|\; s \models g}\). For example, in propositional 
        logic : \( g = At(Robot,RoomB)\wedge Clean(RoomB) \)
    \item \textbf{Goal conditions}: A set of propositions that must be true in the goal state.
        This is the most common approach in classical planning.
    \item \textbf{Utility functions}: In optimization settings, goals may be specified as utility 
        functions to maximize or cost functions to minimize.
\end{itemize}
\begin{definition}{Solution Plan}{Solution Plan}
A \textbf{solution} to a planning problem \( P = (\Sigma,s_0,G) \) is a 
sequence of actions \( \pi = (a_1,a_2,\ldots,a_n) \) such that:
\[
  s_n = \gamma(\gamma(\ldots \gamma(\gamma(s_0,a_1),a_2)\ldots, a_{n-1}),a_n)\in G
\]
That is, executing the sequence of actions starting from the initial state \( s_0 \) results in a goal state.
\end{definition}
We can also write this more compactly using the notation \( \gamma^*(s,\pi) \) to denote
the state reached by executing plan \( \pi \) from state s:
\[
  \gamma^*(s_0,\pi) \in G
\]
\subsubsection{Plan Quality}
Not all solution plans are equally desirable. Common quality metrics include:
\begin{itemize}
    \item \textbf{Plan length}: Number of actions in the plan. Shorter plans are often preferred.
    \item \textbf{Plan cost}: Sum of actions costs. Each action \( a \) may have an associated cost \( c(a) \).
    \item \textbf{Makespan}: Total execution time (relevant in temporal planning).
    \item \textbf{Resource consumption}: Amount of resources used during execution.
\end{itemize}
An \textbf{optimal plan} minimizes the chosen quality metric among all solutions plans.
\section{Representation of Planning Problems}
\subsection{The challenge of explicit representation}
While the state transition system formalism is mathematically elegant, it faces a critical pratical 
challenge: \textbf{the curse of dimensionality}. Real-world planning domains can have exponentially large sate
spaces. For example:
\begin{itemize}
    \item A domain with \( n \) boolean variables has \( 2^n \) possible states 
    \item A domain with \( k \) objects and \( m \) binary relations has up to \( 2^{mk^2} \) states
    \item Enumerating all states and transitions explicitly is infeansible for even moderately-sized problems
\end{itemize}
\begin{example}{Blocks World Complexity}{Blocks World Complexity}
Consider a Blocks World domain with \( n \) blocks. The number of possible 
configurations grows super-exponentially with \( n \). For \( n = 10 \) blocks,
there are more than \( 10^13 \) possible configurations. Explicitly representing the transition 
function would require storing information about trillions of state-action pairs.
\end{example}
This motivates the need for \textbf{compact, structured representation} that exploit regularities 
in the domain to represent large state spaces and transition functions implicitly.
\subsection{Action Schemas: Structured Representation}
The classical approach to compact representation uses \textbf{action schemas} (also called action 
template or operators). An action schema is a prametrized description of a family of related actions.
\begin{definition}{Action Schema}{Action Schema}
An action schema consists of:
\begin{itemize}
    \item \textbf{Name and parameters}: A symbolic name and typed parameters, e.g., \( Move(x,y) \).
    \item \textbf{Preconditions}: A logical formula Pre specifying when the action can be executed
    \item \textbf{Effects}: A logical formula \( Eff \) specifying how the sate changes when the action in executed
\end{itemize} 
\end{definition}
\textbf{Actions} are \textbf{ground instances} of action schemas, obtained by binding the parameters to specific objects in the 
domain. For example, the schema \( Move(x,y) \) might generate ground actions \( Move(RoomA,RoomB) \), \( Move(RoomB,RoomC) \), etc.
\subsubsection{Advantages of Action Schemas}
Action schemas provide severeal benefits:
\begin{itemize}
    \item \textbf{Compactness}: A single schema can represent exponentially many ground actions 
    \item \textbf{Generality}: Schemas capture the general structure of actions independent of specific objects 
    \item \textbf{Scalability}: Adding new objects to the domain automatically generates new appliable actions 
    \item \textbf{Knowledge reuse}: Schemas learned in one domain can be transferred to similar domains.
\end{itemize}
\subsection{State Representation: Propositional logic}
In classical planning, states are typically represented using \textbf{propositional logic}:
\begin{itemize}
    \item A state is a set of \textbf{atoms} (propositional variables) that are true in that state
    \item Atoms not in the set are assumed false (\textbf{closed-world assumption})
    \item Action preconditions and effects are logical formula over these atoms 
\end{itemize}
The representation allows:
\begin{itemize}
    \item Compact encoding of structured states
    \item Efficient reasoning about action applicability and effects 
    \item Use of logical inference techniques for plan generation
\end{itemize}
\subsubsection{State Update Semantics}
\begin{example}{STRIPS Assumption}{STRIPS Assumption}
When an action is executed, the state is updated according to the action's effects:
\begin{itemize}
    \item \textbf{Add list}: Atoms that become true after the action 
    \item \textbf{Delete list}: Atoms that become false after the action
    \item \textbf{Frame axioms}: Atoms not mentioned in the effects remain unchanged
\end{itemize}
The \textbf{STRIPS assumption} (named after the Stanford Research Institute Problem Solver) states
that only atoms explicitly mentioned in the effects changem all other atoms persist unchanged.
This simplifies reasoning about action effects.
\end{example}
\subsection{Example: Blocks World}
Consider the classic \textbf{Blocks World} domain, one of the most studied domain in planning research
\subsubsection{Domain Description}
The Blocks World consists of:
\begin{itemize}
    \item A set of blocks that can be stacked on top of each other 
    \item A table with unlimited space
    \item A robot arm that can pick up and put down blocks
\end{itemize}
\begin{example}{Blocks World State Representation}{Blocks World State Representation}
\textbf{State predicates}:
\begin{itemize}
    \item \( \texttt{On}(x,y) \): Block \( x \) is directly on top of blocks \( y \).
    \item \( \texttt{OnTable}(x) \): Block \( x \) is on the table 
    \item \( \texttt{Clear}(x) \): Block \( x \) has no blocks on top of it 
    \item \( \texttt{Holding}(x) \): The robot arm is holding block \( x \).
    \item \( \texttt{ArmEmpty} \): The robot arm is not holding anything
\end{itemize}
\end{example}
\subsubsection{Action Schemas}
\textbf{PickUp}(x): Pick up block \( x \) from the table 
\begin{align*}
    &\text{Parameters: } x \text{(block)} \\
    &\text{Precondition: } \texttt{Clear}(x) \wedge \texttt{OnTable}(x) \wedge \texttt{ArmEmpty} \\
    &\text{Effect: } \texttt{Holding}(x) \wedge \lnot \texttt{OnTable}(x) \wedge \lnot \texttt{Clear}(x) \wedge \lnot \texttt{ArmEmpty}
\end{align*}
\textbf{PutDown}(x): Put down block \( x \) on the table
\begin{align*}
    &\text{Parameters: } x \text{(block)} \\
    &\text{Preconditions: } \texttt{Holding}(x) \\
    &\text{Effect: } \texttt{OnTable}(x) \wedge \texttt{Clear}(x) \wedge \texttt{ArmEmpty} \wedge \lnot \texttt{Holding}(x)
\end{align*}
\textbf{Stack}(x, y): Stack block \( x \) on top of block \( y \).
\begin{align*}
    &\text{Parameters: } x,y \text{(blocks)} \\
    &\text{Preconditions: } \texttt{Holding}(x) \wedge \texttt{Clear}(x) \\
    &\text{Effect: } \texttt{On}(x,y) \wedge \texttt{Clear}(x) \wedge \texttt{ArmEmpty} \wedge \lnot \texttt{Holding}(x) \wedge \lnot \texttt{Clear}(y)
\end{align*}
\textbf{Unstack}(x, y): Remove block \( x \) from on top of block \( y \).
\begin{align*}
    &\text{Parameters: } x,y \text{(block)} \\
    &\text{Precondition: } \texttt{Holding}(x) \wedge \texttt{Clear}(x) \wedge \texttt{ArmEmpty} \\
    &\text{Effect: } \texttt{Holding}(x) \wedge \texttt{Clear}(y) \wedge \lnot \texttt{On}(x,y) \wedge \lnot \texttt{Clear}(x) \wedge \lnot \texttt{ArmEmpty}
\end{align*}

\subsubsection{Example Problem Instance}
\textbf{Initial state}: Blocks A, B, C are all on the table
\[
  s_0 = \left\{ \texttt{OnTable}(A), \texttt{OnTable}(B), \texttt{OnTable}(C), \texttt{Clear}(A), \texttt{Clear}(B), \texttt{Clear}(C), \texttt{ArmEmpty}\right\}
\]
\textbf{Goal}: Stack the blocks in the order A on B on C 
\[
  G = \left\{ \texttt{On}(A,B), \texttt{On}(B,C), \texttt{OnTable}(C)\right\}
\]
\textbf{Solution plan}:
\begin{enumerate}
    \item \texttt{PickUp}(B)
    \item \texttt{Stack}(B,C)
    \item \texttt{PickUp}(A)
    \item \texttt{Stack}(A,B)
\end{enumerate}

\part{Advanced Planning Representation}
\section{Reduction to Propositional Satisfiability}
\subsection{Motivation}
We continue our discussion of the formalization and representation of planning problems.
A fundamental technique in automated planning is to \textbf{reduce the planning problem to the 
stisfiability problem of propositional logic}. This reduction allows us to leverage powerful SAT
solvers to find plans efficently.
\begin{definition}{Reduction to SAT}{Reduction to SAT}
The key idea is to encode:
\begin{itemize}
    \item The initial state as a propositional formula
    \item The goal conditions as a propositional formula 
    \item The action effects and preconditions as propositional formulas 
    \item The constraint that actions must form a valid sequence
\end{itemize}
If the resulting formula is satisfiable, the satisfying assignment corresponds to a valid plan.
\end{definition}
\subsection{Example Domain: Dock Worker Robots (DWR)}
To illustrate these concepts, we introduce the \textbf{Dock Worker Robots (DWR)} domain, a classic 
planning benchmark that models container logistics at a port.
\subsubsection{Domain description}
The DWR world consists of:
\begin{itemize}
    \item \textbf{Docks}: Three docks \( d_1, d_2, d_3 \) arranged in a triangular configuration 
    \item \textbf{Connectivity}: Docks are connected as follows: \( d_1 \leftrightarrow d_3, d_1 \leftrightarrow d_2, d_2 \leftrightarrow d_3 \).
    \item \textbf{Piles}: Storage locations at docks where containers can be stacked 
    \item \textbf{Containers}: Object that need to be moved between locations 
    \item \textbf{Robots}: Mobile agents that can move containers between piles 
\end{itemize}
\begin{center}
    \begin{tikzpicture}[
        node distance=1.5cm,
        every node/.style={circle, draw, minimum size=1.2cm}
    ]
    \node (d1) {$d_1$};
    \node (d2) [right=of d1] {$d_2$};
    \node (d3) [below=of d1] {$d_3$};

    \draw[<->] (d1) -- (d2);
    \draw[<->] (d1) -- (d3);
    \draw[<->] (d2) -- (d3);
    \end{tikzpicture}
\end{center}

\subsection{Initial State configuration}
The initial state of our example problem is:
\begin{itemize}
    \item \textbf{At dock} \( d_1 \): Pile \( p_1 \) contains two containers: \( c_1 \) (on top) and \( c_2 \) (below)
    \item \textbf{At dock} \( d_2 \):
    \begin{itemize}
        \item Pile \( p_2 \) contains one container: \( c_3 \).
        \item Pile \( p_3 \) is empty 
    \end{itemize}
    \item \textbf{At dock} \( d_3 \): Pile \( p_4 \) is empty
\end{itemize}
\subsection{Formal Representation}
\subsubsection{Objects and Types}
The domain contains objects of different \textbf{sorts} (types). Each object is represented by a constant symbol 
with an associated type:
\vspace{1em}
\noindent
\newline
\textbf{Docks}: 
\[
    \text{Docks } = \left\{ d_1, d_2, d_3 \right\}
\]
where each \( d_i \) is a constant symbol of sort \texttt{Dock}.
\vspace{1em}
\noindent
\newline
\textbf{Containers}:
\[
  \text{Containers } = \left\{ c_1, c_2, c_3 \right\} 
\]
where each \( c_i \) is a constant symbol of sort \texttt{Container}.
\vspace{1em}
\noindent
\newline
\textbf{Robots}:
\[
  \text{Robots } = \left\{ r_1, r_2 \right\}
\]
where each \( r_i \) is a constant symbol of sort \texttt{Robot}.
\vspace{1em}
\noindent
\newline
\textbf{Piles}:
\[
  \text{Piles } = \left\{ p_1, p_2, p_3, p_4 \right\}
\]
where each \( p_i \) is a constant symbol of sort \texttt{Pile}.
\vspace{1em}
\noindent
\newline
The complete set of objects is:
\[
  \text{Objects } = \text{ Docks } \cup \text{ Piles } \cup \text{ Containers } \cup \text{ Robots}
\]
\subsubsection{Static Relations}
Some relations in the domain are \textbf{static} - they do not change during plan exectuion. These can be specified 
once and remain constant throughout.
\vspace{1em}
\noindent
\newline
\textbf{Adjacency relation}: Specifies which docks are directly connected, allowing robots to travel between them.
\[
  \texttt{adjacent} = \left\{ (d_1,d_3), (d_3,d_1), (d_2,d_3), (d_3,d_2), (d_1,d_2), (d_2,d_1) \right\}
\]
This realtion is symmetric: if dock \( d_i \) is adjacent to dock \( d_j \), then \( d_j \) is adjacent to \( d_i \).
\vspace{1em}
\noindent
\newline
\textbf{Pile-Dock association}: Each pile is permanently located at a specific dock.
\begin{align*}
    &\texttt{at}(p_1,d_1)\\
    &\texttt{at}(p_2,d_2)\\
    &\texttt{at}(p_3,d_2)\\
    &\texttt{at}(p_4,d_3)\\
\end{align*}
\subsubsection{Rigid Relations}
The adjacency relation is defined over \( \texttt{Docks} \times \texttt{Docks} \). This is called a \textbf{rigid relation} (also known 
as a \textbf{static relation}) - it doeas not change during plan execution. Actions cannot modify the physical layout of the docks or their connectivity.
\vspace{1em}
\noindent
\newline
Rigid relations are important because:
\begin{itemize}
    \item They reduce the state space by factoring out unchanging aspects 
    \item They can be checked once at planning time rather than during execution 
    \item They simplify action precodinitions by providing fixed background knowledge
\end{itemize}
\subsection{State Representation in DWR}
\subsubsection{What is a State?}
A \textbf{state} is an assignment of values to propositional variables. More precisely, a state can be represented in two equivalent ways:
\begin{enumerate}
    \item \textbf{First-order ground atoms}: Predicates where all arguments are constants (no variables)
    \item \textbf{First-order ground terms}: Terms where all arguments are constants.
\end{enumerate}
\textbf{Important restriction}: We only allow \textbf{flat expression}, meaning expressions with depth 1 (no nesting).
\vspace{1em}
\noindent
\newline
\begin{example}{Forbidden Nesting}{Forbidden Nesting}
Consider the predicate \( \texttt{loc}(x,y) \) meaning "the location of \( x \) is \( y \)". We do \textbf{not} allow 
nested expressions such as:
\[
  \texttt{loc}(r_1,\texttt{loc}(r_2))
\]
This would represent "the locaation of \( r_1 \) is the location of \( r_2 \)", which violates the flatness constraints.
\vspace{1em}
\noindent
\newline
All atoms and terms must be flat - arguments of predicates and functions must be constants, not complex expressions.
\end{example}
\subsubsection{State as Truth Assignment}
A state assign trugth values to ground atoms. Consider describing the initial state \( s_0 \) of out DWR example:
\begin{align*}
    &\texttt{loc}(r_1,d_1)\\
    &\texttt{pos}(p_1,d_1)\\
    &\texttt{pos}(c_2,p_1)\\
\end{align*}
These atoms being listed in the state means they are \textbf{implicitly assigned true}:
\begin{align*}
    &\texttt{loc}(r_1,d_1) \quad \leftarrow \quad \text{true}\\
    &\texttt{pos}(p_1,d_1) \quad \leftarrow \quad \text{true}\\
    &\texttt{pos}(c_2,p_1) \quad \leftarrow \quad \text{true}\\
\end{align*}
Atoms not mentioned in the state are assumed false (closed-world assumption).
\subsection{Two Representation Formalisms}
There are two common approaches to representing planning problems, differing in how they encode state information.
\subsubsection{Classical (Propositional) Representation}
In the \textbf{classical representation}, states are represented using only \textbf{Boolean assignments} to propositional atoms.
Each possible fact about the world is represented as a Boolean variable.
\vspace{1em}
\noindent
\newline
\textbf{Example atoms}:
\begin{itemize}
    \item \( \texttt{loc}(r_1,d_1) \): Robot \( r_1 \) is at dock \( d_1 \) (true/false)
    \item \( \texttt{pos}(c_2,p_1) \): Container \( c_2 \) is at pile \( p_1 \) (true/false)
    \item \( \texttt{empty}(r_1) \): Robot \( r_1 \) is not carrying anything (true/false)
\end{itemize}
\textbf{Characteristics}:
\begin{itemize}
    \item Simple and uniform representation 
    \item Direct correspondence to propositional logic 
    \item Easy to encode as SAT problems 
    \item May require many atoms for complex domains
\end{itemize}
\subsubsection{State Variable Representation}
The \textbf{state variable representation} uses a slightly different syntax based on \textbf{function symbols}. Instead of Boolean atoms, 
we use variables that can take on different values from a finite domain.
\vspace{1em}
\noindent
\newline
\textbf{Example}:
\[
  \texttt{loc}(r_1) = d_1
\]
This means "the location of robot \( r_1 \) is \( d_1 \)". Here, \( \texttt{loc}(r_1) \) is a \textbf{state variable} (function that returns a value), 
and \( d_1 \) is its current value.  
\vspace{1em}
\noindent
\newline
\textbf{Comparison with classical representation}:
\begin{center}
    \begin{tabular}{ |c|c| }
        \hline
        \textbf{Classical} & \textbf{State Variable} \\
        \hline
        \( \texttt{loc}(r_1,d_1) \) = true & \( \texttt{loc}(r_1) = d_1 \) \\
        \( \texttt{loc}(r_1,d_2) \) = false & \\
        \( \texttt{loc}(r_1,d_3) \) = false & \\
        \hline
    \end{tabular}
\end{center}
\textbf{Advantages of state variables}:
\begin{itemize}
    \item More compact: one variable instead of multiple Boolean atoms 
    \item Explicitly represents mutual exclusion (robot can only be at one location)
    \item Natural for domains with multi-value attributes 
    \item Reduces the number of variables in SAT encodings 
\end{itemize}
\textbf{Relationship}: The two representations are \textbf{equivalent in expressive power}. Any problem representable in one 
formalism can be translated to the other. The choice depends on:
\begin{itemize}
    \item The planning algorithm being used 
    \item The structure of the domain 
    \item Efficiency considerations for the specific problem
\end{itemize}
\subsubsection{Example: Complete State Representation}
Let us represent the complete initial state \( s_0 \) of our DWR example in both formalisms.
\vspace{1em}
\noindent
\newline
\textbf{Classical representation} (partial):
\begin{align*}
    &\texttt{loc}(r_1,d_1) = \text{ true}\\
    &\texttt{loc}(r_2,d_2) = \text{ true}\\
    &\texttt{pos}(c_1,p_1) = \text{ true}\\
    &\texttt{pos}(c_2,p_1) = \text{ true}\\
    &\texttt{pos}(c_3,p_2) = \text{ true}\\
    &\texttt{on}(c_1,c_2) = \text{ true}\\
    &\texttt{emtpy}(r_1) = \text{ true}\\
    &\texttt{empty}(r_2) = \text{ true}\\
    &\texttt{top}(c_1,p_1) = \text{ true}\\
    &\texttt{top}(c_3,p_2) = \text{ true}\\
\end{align*}
\textbf{State variable representation}:
\begin{align*}
    &\texttt{loc}(r_1) = d_1 \\
    &\texttt{loc}(r_2) = d_2 \\
    &\texttt{pos}(c_1) = p_1 \\
    &\texttt{pos}(c_2) = p_1 \\
    &\texttt{pos}(c_3) = p_2 \\
    &\texttt{on}(c_1) = c_2 \\
    &\texttt{loaded}(r_1) = \text{ nil}\\
    &\texttt{loaded}(r_2) = \text{ nil}\\
    &\texttt{top}(p_1) = c_1 \\
    &\texttt{top}(p_2) = c_3 \\ 
\end{align*}
Both representations capture the same information but organize it differently. The state variable representation is more compact 
and explicitly enforces constraints (e.g., a robot can only be at one location).
\subsection{Key Insight: States as Boolean Assignments}
The fundamental insight is that regardless of the representation formalism chosen:
\[
  \textbf{A state is fundamentalyy a Boolean assignment}
\]
Wheter we use propositional atoms or state variables, we are ultimately assigning truth values (or equivalently, values from finite domains) to a set of variables.
\section{Action Schemas in DWR}
\subsection{Structure of Action Schemas}
An \textbf{action schema} \( \alpha \) is a template for generating ground actions. It consists of:
\begin{definition}{Action Schema Components}{Action Schema Components}
An action schema \( \alpha \) has the following components:
\begin{itemize}
    \item \textbf{Name and parameters}: \( \alpha \) = \( \texttt{name}(z_1, z_2, \ldots, z_k ) \) where \( z_i \) are typed parameters 
    \item \textbf{Preconditions}: \( \text{Pre}(\alpha) = \left\{ p_1,p_2, \ldots, p_n \right\} \) where each \( p_i \) is a literal 
    \item \textbf{Effects}: \( \text{Eff}(\alpha) = \left\{ q_1,q_2, \ldots, q_m \right\} \) where each \( q_i \) is a literal 
    \item \textbf{Cost}: \( c(\alpha) \in \mathbb{N}_0 \) (typically a non-negative integer)
\end{itemize}
\end{definition}
\textbf{Important constraint}: The parameters \( z_i \) of \( \alpha \) can appear in the precodition literals \( p_j \) and effect 
literals \( q_j \). The arguments in these literals can be either:
\begin{itemize}
    \item \textbf{Constants}: Specific objects from the domain 
    \item \textbf{Parameters}: Variables that will be instantiated 
\end{itemize}
No arbitrary free variables are allowed - all variables must be parameters of the action schema. 
\begin{definition}{Ground Action}{Ground Action}
A \textbf{ground action} \( a \) is obtained by substituting all parameters \( z_i \) in action schema \( \alpha \) with specific 
constants from the domain. The action \( a \) is an \textbf{instance} of \( \alpha \)
\end{definition}
\subsection{Example: Take/Load Action}
Consider the action of a robot taking a container from a pile. We define the action schema:
\[
  \alpha : \texttt{take}(r,c,c',p,d)
\]
\textbf{Parameters}:
\begin{itemize}
    \item \( r \): robot (type \texttt{Robot})
    \item \( c \): container to take (type \texttt{Container})
    \item \( c' \): container below \( c \) (type \( \texttt{Container} \cup \left\{ \texttt{nil} \right\} \))
    \item \( p \): pile (type \texttt{Pile})
    \item \( d \): dock (type \texttt{Dock})
\end{itemize}
\subsubsection{Correcting Static Relations}
Before defining the preconditions, we note that there are \textbf{two rigid (static) relations} in the DWR domain:
\begin{enumerate}
    \item \( \texttt{adjacent}: \texttt{Docks} \times \texttt{Docks} \) - specifies which docks are connected 
    \item \( \texttt{at}: \texttt{Piles} \times \texttt{Docks} \) - specifies which pile is located at which dock
\end{enumerate}
Both relations are rigid because the physical layout of the port does not change during plan execution.
\subsubsection{Preconditions}
For the \texttt{take} action to be executable, the following conditions must hold:
\[
  \mathrm{Pre}(\texttt{take}(r,c,c',p,d)) = 
    \left\{
    \begin{array}{l}
        \texttt{loc}(r) = d \\
        \texttt{at}(p,d) \\
        \texttt{pos}(c) = p \\
        \texttt{on}(c) = c' \\
        \texttt{top}(p) = c \\
        \texttt{cargo}(r) = \texttt{nil}
    \end{array}
    \right\}
\]
\textbf{Interpretation}:
\begin{itemize}
    \item The robot \( r \) is at dock \( d \)   
    \item The pile \( p \) is located at dock \( d \) (static relation)
    \item Container \( c \) is positioned at pile \( p \)    
    \item Container \( c \) is on top of container \( c' \) (or \( c' \) is \texttt{nil} if \( c \) is at the bottom)
    \item Container \( c \) is the top container of pile \( p \)    
    \item The robot \( r \) is not carrying anything
\end{itemize}
\subsubsection{Effects}
When the \texttt{take} action is executed, the following changes occur:
\[
  \mathrm{Eff}(\texttt{take}(r,c,c',p,d)) = 
    \left\{
    \begin{array}{l}
        \texttt{cargo}(r) = c \\
        \texttt{top}(p) = c' \\
        \texttt{pos}(c) = r
    \end{array}
    \right\}
\]
\textbf{Interpretation}:
\begin{itemize}
    \item The robot \( r \) is now carrying container \( c \)    
    \item The top of pile \( p \) is now \( c' \) (the container that was below $c$)
    \item The position of container \( c \) is now the robot \( r \)
\end{itemize}
\subsection{Action Applicability}
\begin{definition}{Action Applicability}{Action Applicability}
A ground action \( a \) is \textbf{applicable} (or \textbf{executable}) in state \( s \) id all preconditions of \( a \) are 
satisfied in \( s \):
\[
  a \text{ applicable in } s \quad \Leftrightarrow \quad s \models \mathrm{Pre}(a)
\]
\end{definition}
\subsubsection{Examples}
Consider two ground actions:
\vspace{1em}
\noindent
\newline
\textbf{Example 1}: \( a_1 = \texttt{take}(r_1,c_1,c_2,p_1,d_1) \)  
\vspace{1em}
\noindent
\newline
In the initial state \( s_0 \) where:
\begin{itemize}
    \item $\texttt{loc}(r_1) = d_1$
    \item $\texttt{at}(p_1,d_1)$
    \item $\texttt{pos}(c_1) = p_1$
    \item $\texttt{on}(c_1) = c_2$
    \item $\texttt{top}(c_1) = c_1$
    \item $\texttt{cargo}(r_1) = \texttt{nil}$
\end{itemize}
All preconditions are satisfied, so \( a_1 \) is \textbf{executable/applicable} in \( s_0 \).
\vspace{1em}
\noindent
\newline
\textbf{Example 2}: \( a_2 = \texttt{take}(r_2,c_3,\texttt{nil},p_4,d_3) \)    
\vspace{1em}
\noindent
\newline
This action requires:
\begin{itemize}
    \item $\texttt{loc}(r_2) = d_3$
    \item $\texttt{at}(p_4,d_3)$
    \item $\texttt{pos}(c_3) = p_4$
    \item $\texttt{top}(p_4) = c_3$
\end{itemize}
However, in \( s_0 \), we have \( \texttt{pos}(c_3) = p_2 \) (not $p_4$) and $\texttt{loc}(r_2) = d_2$ (not $d_3$). Therefore, 
$a_2$ is \textbf{not executable/not applicable} in $s_0$.
\section{State Transition Function}
\subsection{Formal Definition}
The \textbf{state transition function} \( \gamma \) defines how actions transform states:
\[
  \gamma : \text{ State } \times \text{ Action } \rightarrow \text{ State}
\]
where \( \gamma(s,a) \) computes the successor state obtained by executing ground action \( a \) in state \( s \).
\subsection{Computing the Successor State}
Given a state \( s \) and a ground action \( a \) (which is a ground instance of an action schema \( \alpha \)), the 
successor state \( \gamma(s,a) = s'\) is defined if and only if:
\begin{enumerate}
    \item \( s \models p \) for all \( p \in \mathrm{Pre}(a) \) (all preconditions are satisfied)
    \item \( s' \) is consistent (the resulting state contains no contradictions)
\end{enumerate}
The successor state \( s' \) is computed as:
\[
  s' = \left\{ q \;|\; q \in \mathrm{Eff}(a) \right\} \cup \left\{ q \;|\; s \models q, \lnot q \notin \mathrm{Eff}(a), q \notin \mathrm{Eff}(a) \right\}
\]
\textbf{Interpretation}: The new state consists of:
\begin{itemize}
    \item All literals in the effects of $a$
    \item All literals from $s$ that are not affected by $a$ (neither the literal nor its negation appears in the effects)
\end{itemize}
\subsection{Goal States}
While the initial state is a specific assignment, goals are typically specified more flexibly.
\begin{definition}{Goal Formula}{Goal Formula}
A \textbf{goal formula} \( G \) is a conjunction of literals specifying desired properties of goal states.
\end{definition}
\begin{definition}{Goal States}{Goal States}
The set of \textbf{goal states} is:
\[
  S_g = \left\{ s \;|\; s \models G \right\}
\]
That is, all assignments that satisfy the goal formula \( G \).
\end{definition}
\subsection{Comparing Representations: Simplified DWR example}
Let us compare the classical and state variable representations using a simplified version of the DWR domain with three actions schema 
\subsubsection{Classical Representation}
In the classical representation, we use propositional atoms to represent facts:
\vspace{1em}
\noindent
\newline
\textbf{Action 1}: \( \texttt{take}(r,c,l) \) - robot \( r \) takes container \( c \) at location $l$
\begin{align*}
    &\mathrm{Pre}: & \texttt{loc}(r,l) \wedge \texttt{pos}(c,l) \wedge \lnot \texttt{loaded}(r) \\
    &\mathrm{Eff}: & \texttt{pos}(c,r) \wedge \lnot \texttt{pos}(c,l) \wedge \texttt{loaded}(r)
\end{align*}
\textbf{Action 2}: \( \texttt{move}(r,l,m) \) - robot \( r \) moves from location \( l \) to location $m$ 
\begin{align*}
    &\mathrm{Pre}: & \texttt{loc}(r,l)\\
    &\mathrm{Eff}: & \texttt{loc}(r,m)\wedge \lnot \texttt{loc}(r,l)\\
\end{align*}
\textbf{Action 3}: \( \texttt{put}(r,c,l) \) - robot \( r \) puts down container \( c \) at location \( l \)    
\begin{align*}
    &\mathrm{Pre}: & \texttt{loc}(r,l) \wedge \texttt{pos}(c,r)\\
    &\mathrm{Eff}: & \texttt{pos}(c,l) \wedge \lnot \texttt{pos}(c,r) \wedge \lnot \texttt{loaded}(r)\\
\end{align*}
\subsubsection{State Variable Representation}
In the state variable representation, we use functional notation with assignments:
\vspace{1em}
\noindent
\newline
\textbf{Action 1}: \( \texttt{take}(r,c,l) \)   
\begin{align*}
    &\mathrm{Pre}: & \texttt{loc}(r) = l \wedge \texttt{pos}(c) = l \wedge \lnot \texttt{loaded}(r) \\
    &\mathrm{Eff}: & \texttt{pos}(c) \leftarrow r \wedge \texttt{loaded}(r) \\
\end{align*}
\textbf{Action 2}: \( \texttt{move}(r,l,m) \)    
\begin{align*}
    &\mathrm{Pre}: & \texttt{loc}(r) = l \\
    &\mathrm{Eff}: & \texttt{loc}(r) \leftarrow m 
\end{align*}
\textbf{Action 3}: \( \texttt{put}(r,c,l) \)   
\begin{align*}
    &\mathrm{Pre}: & \texttt{loc}(r) = l \wedge \texttt{pos}(c) = r\\
    &\mathrm{Eff}: & \texttt{pos}(c) \leftarrow l \wedge \lnot \texttt{loaded}(r)
\end{align*}
\subsubsection{Key Differences}
\begin{itemize}
    \item \textbf{Classical}: Requires explicit deletion of old values (e.g., \( \lnot \texttt{loc}(r,l) \) when moving) 
    \item \textbf{State Variable}: Implicit deletion through reassignment (e.g., \( \texttt{loc}(r) \leftarrow m \) automatically removes the old location) 
    \item \textbf{Compactness}: State variables are more concise for multi-valued attributes
\end{itemize}
\section{Planning as Satisfiability (SAT)}
\subsection{Recap: Core Concepts}
Before introducting the SAT encoding, let us recap the fundamental concepts:
\begin{itemize}
    \item \textbf{State}: A Boolean assignment to propositional (ground) literals
    \item \textbf{Action}: Characterized by:
    \begin{itemize}
        \item \textbf{Preconditions}: A set (or conjunction) of literals 
        \item \textbf{Effects}: A set (or conjuction) of literals
    \end{itemize}
\end{itemize}
\subsection{Temproal Instantiation of Variables}
To encode planning problems as SAT, we introduce a temporal dimensio to our propositional variables.
\begin{definition}{Temporal Instantiation}{Temporal Instantiation}
Let \( \mathcal{X} \) be the set of propositional variables describing the domain. The \textbf{temporal instantiation} 
of \( \mathcal{X} \) at time \( t \) is:
\[
  \mathcal{X}\symbol{64}t = \left\{ x\symbol{64}t \;|\; x \in \mathcal{X} \right\} \quad \text{ for } t \geq 0
\] 
where \( x\symbol{64}t \) is a propositional variable representing the truth value of \( x \) at time step \( t \).
\end{definition}
\textbf{Intuition}: Since a plan is a sequence of actions \( (a_0,a_1,a_2, \ldots, a_{n - 1}) \) of length \( n \), 
we need to track truth value of each proposition at each time step from \( 0 \) to \( n \).
\subsection{The Bounded Planning Problem}
\begin{definition}{Bounded Planning Problem}{Bounded Planning Problem}
Given a planning problem \( P = \langle \Sigma, I, G \rangle \), the \textbf{bounded planning problem} asks:
\vspace{0.5em}
\noindent
\newline
\textit{Does there exists a plan solving P with length exactly n?}
\end{definition}
This is a decision problem parametrized by the plan length \( n \).
\subsection{SAT Encoding}
The bounded planning problem can be encoded as a propositional satisfiability problem.
\begin{theorem}{Planning as SAT}{Planning as SAT}
\textit{For a planning problem P and bound n, there exists a propositional formula} \( \Phi_{n} \) \textit{such that}:
\[
  \Phi_{n} \textit{is satisfiable} \quad \Leftrightarrow \quad \textit{there exists a plan of length n for P}
\]
\end{theorem}
The formula \( \Phi_{n} \) encodes:
\begin{enumerate}
    \item \textbf{Initial state}: \( I\symbol{64}0 \) (the initial state holds at time 0)
    \item \textbf{Goal state}: \( G\symbol{64}n \) (the goal holds at time $n$)
    \item \textbf{Action preconditions and effects}: For each time step \( t \in \left\{0,1, \ldots, n - 1 \right\} \), exactly one action is executed and its effects are correctly applied
    \item \textbf{Frame Axioms}: Variables not affected by actions remain unchanged
\end{enumerate}
\subsection{Solving Planning via Iterative Deepening}
Since we don't know the optimal plan length in advance, we use \textbf{iterative deepening}:
\begin{algorithm}[H]
    \caption{Planning via SAT with Iterative Deepening}
    \begin{algorithmic}[1]
      \State \( n \leftarrow 1 \)
      \While{true}
        \State Construct formula \( \Phi_{n} \)
        \If{\( \Phi_{n} \) is satisfiable}
          \State Extract plan from satisfying assignment
          \State \Return plan
        \EndIf
        \State \( n \leftarrow n + 1 \)
      \EndWhile
    \end{algorithmic}
\end{algorithm}
\textbf{Process}:
\begin{enumerate}
    \item Try \( n = 1 \): Check if \( \Phi_{1} \) is satisfiable 
    \item If not, try \( n = 2 \): Check if \( \Phi_2 \) is satisfiable 
    \item If not, try \( n = 3 \): Check if \( \Phi_3 \) is satisfiable 
    \item Continue until a satisfiable formula is found 
\end{enumerate}
\textbf{Advantages}:
\begin{itemize}
    \item Finds shortest plan (optimal in number of steps)
    \item Leverage highly opimized SAT solvers 
    \item Scales well for many planning domains 
\end{itemize}
\textbf{Challenges}:
\begin{itemize}
    \item Formula size grows with $n$ 
    \item May need to try many values of \( n \) before finding a solution 
    \item Encoding frame axioms can be expensive 
\end{itemize}
\part{SAT Solvers and Planning}
\section{The Boolean Satisfiability Problem}
\subsection{Definition and Scope}
The \textbf{Boolean Satisfiability Problem (SAT)} is one of the most fundamental problems in computer science and forms 
the foundation for many automated reasoning techniques, including planning 
\begin{definition}{SAT Problem}{SAT Problem}
The SAT problem asks: Given a propositional formula, does there exists an interpretation (truth assignment) that makes 
the formula true? Formally, given a seto of clauses \( S \) in propositional logic, determine whether \( S \) is satisfiable.
\end{definition}
\subsection{Propositional Logic Basics}
To understand SAT, we must first establish the basic elements of propositional logic:
\begin{itemize}
    \item \textbf{Propositional variables}: Let \( \mathcal{X} \) be a finite set of propositional variables, where each variable can assume one of two values: \texttt{true}(1) or \texttt{false}(0)
    \item \textbf{Atom (Atomic formula)}: A single propositional variable \( x \in \mathcal{X} \)    
    \item \textbf{Literal}: Either an atom \( x \) or its negation \( \lnot x \)     
    \item \textbf{Clause}: A disjunction of literals: \( L_1 \vee L_2 \vee, \ldots, \vee L_n \) where each \( L_i \) is a literal
\end{itemize}
\subsection{Interpretation and Truth Assignments}
\begin{definition}{Interpretation}{Interpretation}
An \textbf{interpretation} (or \textbf{truth assignment}) is a function:
\[
  v: \mathcal{X} \rightarrow \left\{ 0,1 \right\}
\]
that assign a Boolean value to each propositional variable
\end{definition}
\textbf{Number of interpretations}: If \( \;|\;\mathcal{X}\;|\; = n \) (i.e., we have $n$ propositional variables), then there are exactly 
\( 2^n \) possible interpretations.
\subsection{Satisfiability of Clauses}
\begin{definition}{Satisfaction}{Satisfaction}
A clause \( C = L_1 \vee L_2 \vee \ldots \vee L_n \) is \textbf{satisfied} by an interpretation \( v \) if there 
exists at least one literal \( L_i \) such that \( v(L_i) = \text{ true} \) for some \( i \in \left\{ 1,2, \ldots, n \right\} \).
\vspace{1em}
\noindent
\newline
We write: \( v \models C \). 
\vspace{1em}
\noindent
\newline
A set of clauses \( S = \left\{ C_1, C_2, \ldots, C_m \right\} \) is \textbf{satisfiable} if there exists an interpretation 
\( v \) that satisfies all clauses in \( S \):
\[
  v \models S \quad \Leftrightarrow \quad v \models C_i \text{ for all } i \in \left\{1,2, \ldots, m \right\} 
\]
\end{definition}
\subsection{SAT Decision Procedure}
The SAT problem can be viewed as a decision procedure with the following input-output specification:
\fbox{%
    $\displaystyle 
    S\ \text{(set of clauses)}
    \;\xrightarrow{\;\text{SAT Solver}\;}\;
    \left\{
        \begin{array}{l}
            \textbf{YES (Satisfiable)}\\[2pt]
            \exists v : v \models S\\[10pt]
            \textbf{No (Unsatisfiable)}\\[2pt]
            \text{Proof: } S \vdash \square
        \end{array}
    \right.
    $%
}
\textbf{Output cases}:
\begin{itemize}
    \item \textbf{YES (Satisfiable)}: There exists an interpretation \( v \) such that \( v \models S \). The SAT solver typically returns this satisfying assignment.
    \item \textbf{NO (Unsatisfiable)}: No interpretation satisfies \( S \). The solver may provide a proof of unsatisfiability, showing that \( S \) derives a contradiction (denoted \( S \vdash \square \), where \( \square \) represents the empty clause).
\end{itemize}
\subsection{The Empty Clause}
\begin{definition}{Empty Clause}{Empty Clause}
The \textbf{empty clause}, denoted \( \square \) (or \( \square \)), is a clause with no literals. It is \textbf{false in all interpretation}
and therefore represents a contradiction.
\end{definition}
\textbf{Significance}: If a SAT solver derives the empty clause from the input formula \( S \), this constitutes a proof that \( S \) is unsatisfiable, 
since deriving a contradiction means no interpretation can satisfy all clauses.
\section{Planning Problem Formalization}
We formalized the planning problem with:
\[
  P = \langle \Sigma, s_o, S_g \rangle
\]
where \( \Sigma = \langle S, A, \gamma \rangle \)
\begin{itemize}
    \item \( s_0 \) is the initial state 
    \item \( S_g \) is the set of goal states: \( S_g = \left\{ s \;|\; s \models G \right\} \)    
    \item \( G \) is the goal formula (a conjuction of atoms) 
\end{itemize}
\subsection{States and Interpretations}
States correspond to interpretations in propositional logic:
\begin{itemize}
    \item \( \texttt{at}(r_1, l_1) \): ground positive literal 
    \item \( \lnot \texttt{loaded}(r_1) \): ground negative literal
\end{itemize}
\textbf{Abstraction}: From ground first-order literals to propositional literals. With finitely many ground atoms (literals) in the planning 
problem, we obtain finitely many variables \( \;|\; X \;|\; = n \).
\section{Bounded Planning Problem}
\begin{definition}{Bounded Planning Problem}{Bounded Planning Problem}
Does there exists a plan \( \pi \) that solves planning problem \( P \) and has length \( k \) (i.e., $\pi = k$, where $\pi$ is a sequence of $k$ actions)?
\end{definition}
\subsection{Iterative Deepening on Plan Length}
We use iterative deepening to find the shortest plan by increasing the bound \( k \):
\fbox{%
    $\displaystyle
    P\      
    \; \xrightarrow{\;\text{translation}\;}\;
    \Phi\       
    \; \xrightarrow{\;\text{SAT}\;}\;
    \left\{      
        \begin{array}{l}
            \textbf{YES}: \text{ interpretation } v \models \Phi 
            \text{ encodes plan } \pi \text{ of length } k\\[10pt]
            \textbf{No}: \text{ no plan of length } k \text{ solves } P
        \end{array}
    \right.
    $%
}
The process iterates \( k := k + 1 \) until a solution is found. 
\vspace{1em}
\noindent
\newline
\textbf{Key Insight}: The SAT encoding \( \Phi \) includes:
\begin{itemize}
    \item Initial state constraints at time 0 
    \item Goal constraints at time \( k \)     
    \item Action precodinitions and effects for each time step 
    \item Frame axioms ensuring persistence of unchanged literals
\end{itemize}
\subsection{Temporal Encoding of Literals}
A key aspect of the SAT encoding for planning is the \textbf{temporal instantiation} of literals. Since a plan of length \( k \) requires tracking 
state changes over time, each literal is indexed by time steps.
\begin{itemize}
    \item \textbf{Plan length} \( k \) we have \( 0 \leq i \leq k - 1 \) time steps (transitions)
    \item Each literals \( f \) becomes \( f\symbol{64}t \) where \( t \) ranges from \( 0 \) to \( k \)     
    \item State \( i \) is the state after executing \( i \) actions (before the \( i \)-th action)
\end{itemize}
\subsubsection{Temporal Instantiation}
Consider the literal \( \texttt{at}(r_1, l_1) \) ("robot \( r_1 \) is at location \( l_1 \)"):
\[
  \texttt{at}(r_1, l_1) \xrightarrow{\text{temporal encoding}} \texttt{at}(r_1, l_1, i)
\]
Similarly, the negative literal \( \lnot \texttt{loaded}(r_1) \) becomes \( \lnot \texttt{loaded}(r_1, i) \), meaning "robot" \( r_1 \) is not 
loaded at stage \( i \)".
\begin{definition}{State Transition}{State Transition}
Step \( i \) (where $0 \leq i \leq k - 1$) takes us from state \( i \) to state \( i + 1 \) by executing one action.
\end{definition}
\subsection{Encoding the Initial State}
The initial state \( s_0 \) is encoded as a conjuction of constraints at time \( t = 0 \):
\[
  \Phi_{s_0} = \bigwedge_{f \in s_0} f\symbol{64}0 \quad \wedge \quad \bigwedge_{f \notin s_0} \lnot f\symbol{64}0
\]
Where \( f \) represents ground literals.
\subsubsection{From Literals to Propositional Variables}
Each temporally indexed literal corresponds to a propositional variable:
\begin{itemize}
    \item \( \texttt{at}(r_1,l_1,0) \) correspond to variable \( x_{\texttt{at}(r_1,l_1)}\symbol{64}0 \)
    \item \( \lnot \texttt{loaded}(r_1,0) \) correspond to \( \lnot x_{\texttt{loaded}(r_1)}\symbol{64}0 \)
\end{itemize}
\subsubsection{Example: Complete Initial State Encoding}
Consider an initial state containing \( \texttt{at}(r_1,l_1) \) but not \( \texttt{laoded}(r_1) \):
\begin{align*}
    & \text{Literals in } s_0: & \texttt{at}(r_1,l_1) & \wedge \lnot \texttt{loaded}(r_1)\wedge \ldots \\
    & \text{Temporal encoding }: & \texttt{at}(r_1,l_1,0) & \wedge \lnot \texttt{loaded}(r_1,0) \wedge \ldots \\
    & \text{SAT variables }: & x_{\texttt{at}(r_1,l_1)}\symbol{64}0 & \wedge \lnot x_{\texttt{loaded}(r_1)}\symbol{64}0\wedge \ldots \\
\end{align*}
For literals not in the initial state (e.g., $\texttt{at}(r_1,l_2), \texttt{at}(r_1,l_3$)), these are assigned false:
\[
  \lnot \texttt{at}(r_1,l_2,0), \quad \lnot \texttt{at}(r_1,l_3,0)
\]
\textbf{Alternative notation}: The initial state encoding can also be written using set notation:
\[
  \Phi_{s_0} = \bigwedge_{x \in S_0^+} x_0 \wedge \bigwedge_{x \in S_0^-} \lnot x_0
\]
where \( S_0^+ \) is the set of positive literals in \( s_0 \) and \( S_0^- \) is the set of negative literals in \( s_0 \).
\subsection{Encoding the Goal State}
The goal conditions must be satisfied at the final time step \( k \).
\begin{definition}{Goal Encoding}{Goal Encoding}
The goal states are encoded at time \( k \) as:
\[
  \Phi_G = \bigwedge_{x \in G^+} x_k \wedge \bigwedge_{x \in G^-} \lnot x_k
\]
where \( G^+ \) is the set of positive literals in the goal \( G \) and \( G^- \) is the set of negative literals in \( G \).
\end{definition}
\textbf{Interpretation}: This formula ensures that all positive goal conditions are true and all negative goal conditions are false 
i nthe final state after executing \( k \) actions.
\subsection{Encoding Actions}
The planning problem is defined as \( \Sigma = \langle S, A, \gamma \rangle \) where \( A \) is the set of all ground instances of 
the action schemas.
\subsubsection{Action Preconditions}
For each action \( a \in A \) at each time step \( i \) (where $0 \leq i \leq k - 1$), we encode the constraint that 
if the action is executed, its preconditions must hold:
\begin{definition}{Action Precondition Encoding}{Action Precondition Encoding}
\[
  \forall a \in A, \forall i \in \left\{0, \ldots, k - 1 \right\}: \quad a_i \rightarrow \bigwedge_{p \in \text{prec}(a)}p_i
\]
where \( \text{prec}(a) \) denotes the set of preconditions of action \( a \).
\end{definition}
\textbf{Interpretation}: This formula states that if action \( a \) is executed at time \( i \), then all of its preconditions 
must be satisfied at time \( i \). The implication can be rewritten in CNF as:
\[
  \lnot a_i \vee \bigwedge_{p \in \text{prec}(a)} p_i
\]
\subsubsection{Mutual Exclusion of Actions}
At each time step, at most one action can be executed. This constraints is encoded as follows:
\begin{definition}{Action Mutual Exclusion}{Action Mutual Exclusion}
For all pairs of distinct actions $(a,b) \in A \times A$ where \( a \neq b \), and for all time steps \( i \) where \( 0 \leq i \leq k - 1 \):
\[
  \lnot a_i \vee \lnot b_i
\]
\end{definition}
This formula ensures that actions \( a \) and \( b \) cannot both be executed at the same time step \( i \).
\vspace{1em}
\noindent 
\newline
\textbf{Equivalent formulas}:
\begin{itemize}
    \item \( a_i \rightarrow \lnot b_i \): If action \( a \) is executed at time \( i \), then action \( b \) is not executed 
    \item \( b_i \rightarrow \lnot a_i \): If action \( b \) is executed at time \( i \), the action \( a \) is not executed
\end{itemize} 
\subsection{Complete SAT Encoding}
The complete SAT encoding \( \Phi \) for a bounded planning problem with plan length \( k \) is the conjunction of the following sets of propositional fromulas:
\[
  \Phi = \Phi_{s_0} \wedge \Phi_G \wedge \Phi_{\text{actions}} \wedge \Phi_{\text{mutex}} \wedge \Phi_{\text{frame}}
\]
where: 
\begin{itemize}
    \item \( \Phi_{s_0} \): Initial state constraints 
    \item \( \Phi_G \): Goal state constraints 
    \item \( \Phi_{\text{actions}} \): Action precondition and effect constraints 
    \item \( \Phi_{\text{mutex}} \): Mutual exclusion constraints between actions 
    \item \( \Phi_{\text{frame}} \): Frame axioms (persistence of unchanged literals)
\end{itemize}
\subsection{Plan Extraction from SAT Model}
Once the SAT solver returns a satisfying assignment, we can extract the corresponding plan.
\begin{definition}{Plan Extraction}{Plan Extraction}
Let \( v \) be a satisfying interpretation (model) returned by the SAT solver, represented as a set of literals. The plan \( \pi \) is extracted as follows:
\[
  \pi = \langle a_0, a_1, \ldots, a_{k - 1} \rangle
\]    
where \( a_i \) is the unique action such that \( v \models a_i \) (i.e., $a_i$ is assigned \texttt{true} in the model $v$) for each time 
step \( i \in \left\{ 0,1,\ldots,k-1 \right\} \).
\end{definition}
\textbf{Extraction procedure}:
\begin{enumerate}
    \item Filter the model \( v \) to obtain only the action literals: \( \left\{ a_i \;|\; a \in A, 0 \leq i \leq k - 1, v \models a_i \right\} \)     
    \item For each time step \( i \), identify the action \( a \) such that \( a_i \) is true in \( v \)      
    \item Construct the sequence \( \pi = \langle a_0, a_1, \ldots, a_{k - 1} \rangle \)    
\end{enumerate}
\textit{Remark}. Due to the mutual exclusion constraints, at most one action is true at each time step. If no action is true at some time step 
\( i \), this correspond to a "no-op" (no operation) action.
\part{SAT Encoding Examples}
\section{Detailed SAT Encoding Example}
\subsection{Simple Robot Navigation Domain}
To illustrate the complete SAT encoding process, we present a simple example with a single 
robot ad two locations.
\vspace{1em}
\noindent
\newline
\textbf{Example} (Robot Navigation). Consider a minimal planning domain with the following components:
\vspace{1em}
\noindent
\newline
\textbf{Objects}:
\begin{itemize}
    \item One robot: \( r_1 \)    
    \item Two locations: \( l_1, l_2 \) 
\end{itemize}
\textbf{Action schema}: \( \texttt{move}(r,l,l') \)    
\begin{align*}
    &\text{Parameters:} & r(\text{robot}), l(\text{robot}),l(\text{source location}),l'(\text{destination location})\\
    &\text{Preconditions:} & \texttt{at}(r,l)\\
    &\text{Effect:} & \texttt{at}(r,l') \wedge \lnot \texttt{at}(r,l)
\end{align*}
\subsubsection{Ground actions}
With one robot and two locations, we obtain two ground actions:
\begin{enumerate}
    \item \( a_1 = \texttt{move}(r_1,l_1,l_2) \): Move robot \( r_1 \) from location \( l_1 \) to location \( l_2 \)    
    \item \( a_2 = \texttt{move}(r_1,l_2,l_1) \): Move robot \( r_1 \) from location \( l_2 \) to location \( l_1 \)
\end{enumerate}
\subsubsection{Problem Instance}
\textbf{Initial state} (at time step $t = 0$):
\[
  s_0 = \left\{\texttt{at}(r_1,l_1) \right\}
\]
The robot starts at location \( l_1 \). This is encoded temporally as:
\[
  \texttt{at}(r_1,l_1)\symbol{64}0
\]
\textbf{Goal state} (at time step $t = 1$):
\[
  G = \left\{ \texttt{at}(r_1,l_2) \right\}
\]
The goal is for the robot to be at location \( l_2 \) after one action. This is encoded as:
\[
  \texttt{at}(r_1,l_2)\symbol{64}1
\]
\subsection{Complete SAT Encoding}
For this problem with plan length $k = 1$, we construct the following propositional formulas:
\subsubsection{Initial State COnstraints}
\[
  \Phi_{s_0} = \texttt{at}(r_1,l_1)\symbol{64}0 \wedge \lnot \texttt{at}(r_1,l_2)\symbol{64}0
\]
\subsubsection{Goal state Constraints}
\[
  \Phi_G = \texttt{at}(r_1,l_2)\symbol{64}1
\]
\subsubsection{Action Preconditions and Effect}
For action \( a_1 = \texttt{move}(r_1,l_1,l_2) \) at time step 0:
\[
  \texttt{move}(r_1,l_1,l_2)\symbol{64}0 \rightarrow (\texttt{at}(r_1,l_1)\symbol{64}0 \wedge \texttt{at}(r_1,l_2)\symbol{64}1 \wedge \lnot \texttt{at}(r_1,l_1)\symbol{64}0)
\]
For action \( a_2 = \texttt{move}(r_1,l_2,l_1) \) at time step 0:
\[
  \texttt{move}(r_1,l_2,l_1)\symbol{64}0 \rightarrow (\texttt{at}(r_1,l_2)\symbol{64}0 \wedge \texttt{at}(r_1,l_1)\symbol{64}1 \wedge \lnot \texttt{at}(r_1,l_2)\symbol{64}1)
\]
\subsubsection{Explanatory Frame Axioms}
The \textbf{explanatory frame axioms} (also called \textbf{explanation axioms}) state that if a literal changes its truth
value between time steps, then some action causing that change must have been executed.
\vspace{1em}
\noindent
\newline
\textbf{If} \( \texttt{at}(r_1,l_1) \) \textbf{becomes false}:
\[
  \texttt{at}(r_1,l_1)\symbol{64}0 \wedge \lnot \texttt{at}(r_1,l_1)\symbol{64}1 \rightarrow \texttt{move}(r_1,l_1,l_2)\symbol{64}0
\]
\textbf{If} \( \texttt{at}(r_1,l_2) \) \textbf{becomes false}:
\[
  \texttt{at}(r_1,l_2)\symbol{64}0 \wedge \lnot \texttt{at}(r_1,l_2)\symbol{64}1 \rightarrow \texttt{move}(r_1,l_2,l_1)\symbol{64}0
\]
\textbf{If} \( \texttt{at}(r_1,l_1) \) \textbf{becomes true}:
\[
  \lnot \texttt{at}(r_1,l_1)\symbol{64}0 \wedge \texttt{at}(r_1,l_1)\symbol{64}1 \rightarrow \texttt{move}(r_1,l_2,l_1)\symbol{64}0
\]
\textbf{If} \( \texttt{at}(r_1,l_2) \) \textbf{becomes true}:
\[
  \lnot \texttt{at}(r_1,l_2)\symbol{64}0 \wedge \texttt{at}(r_1,l_2)\symbol{64}1 \rightarrow \texttt{move}(r_1,l_1,l_2)\symbol{64}0
\]
\subsection{Solution Extraction}
Given the constraints, the SAT solver will find a satisfying assignment that includes:
\[
  \texttt{move}(r_1,l_1,l_2)\symbol{64}0 = \text{ true}
\]
The extracted plan is:
\[
  \pi = \langle \texttt{move}(r_1,l_1,l_2) \rangle
\]
This plan has length 1 and achieves the goal of moving the robot from \( l_1 \) to \( l_2 \).
\section{Conversion to Conjunctive Normal Form}
\subsection{CNF definition and Structure}
To use a SAT solver, we must convert all formulas to \textbf{Conjunctive Normal Form (CNF)}.
\begin{definition}{Conjunctive Normal Form (CNF)}{Conjunctive Normal Form (CNF)}
A \textbf{clause} is a disjunction of literals: \( L_1 \vee L_2 \vee \ldots \vee L_k \) where 
each \( L_i \) is a literal (an atom or its negation)
\vspace{1em}
\noindent
\newline
A formula is in \textbf{Conjunctive Normal Form} if it is a conjunction of clauses:
\[
  \text{CNF } = (L_1^1 \vee L_2^1 \vee \ldots \vee L_k^1) \wedge (L_1^2 \vee L_2^2 \vee \ldots \vee L_k^2) \wedge \ldots \wedge (L_1^m \vee L_2^m \vee \ldots \vee L_k^m)
\]
\end{definition}
A set of clauses is logically equivalent to a conjunction of those clauses, which is precisely the CNF representation 
required by SAT solvers.
\subsection{Standard CNF Transformation}
THe standard transformation to CNF involves applying logical equivalences to eliminate implications and biconditionals,
push negations inward, and distribute disjunctions over conjunctions.
\subsubsection{Logical Equivalences}
The following logical equivalences are used in this transformation:
\vspace{1em}
\noindent
\newline
\textbf{Biconditional elimination}:
\[
  F \leftrightarrow G \equiv (F \rightarrow G) \wedge (G \rightarrow F)
\]
\vspace{1em}
\noindent
\newline
\textbf{Implication elimination}:
\[
  F \rightarrow G \equiv \lnot F \wedge G 
\]
\vspace{1em}
\noindent
\newline
\textbf{De Morgan's laws}:
\begin{align*}
    & \lnot(F \vee G) & \equiv & \lnot F \wedge \lnot G \\
    & \lnot(F \wedge G) & \equiv & \lnot F \vee \lnot G \\
\end{align*}
\subsubsection{Negation Normal Form (NNF)}
After applying the above equivalences, we obtain a formula in \textbf{Negation Normal Form (NNF)}, where negations appear only 
directly in front of atoms (not in front of complex formulas).
\subsubsection{Distribution to Obtain CNF}
To convert NNF to CNF, we apply the \textbf{distributive law}:
\[
  F \vee (G \wedge H) \equiv (F \vee G) \wedge (F \vee H)
\]
This distributes disjunctions over conjunctions, eventually producing a conjunction of disjunctions (CNF).
\subsection{The Exponential Blowup Problem}
\textbf{Critical issue}: The standard CNF transformation cna cause a \textbf{combinatorial explosion} in formula size.
\begin{theorem}{Exponential Blowup Problem}{Exponential Blowup Problem}
\textit{Observation} (Exponential Blowup). In the wrost case, the standard transformation into CNF produces a formula 
whose size is \textbf{exponential} in the size of the original formula:
\newline
\( O(2^n) \) where \( n \) is the size of the input formula.
\end{theorem}
This exponential blowup makes the standard transformation \textbf{impratical} fro large planning problems, as 
the resulting CNF formula becomes too large to handle efficiently.
\subsection{Equisatisfiable CNF (ECNF)}
To avoid the exponential blowup, we use a different approach that produces an \textbf{Equisatisfiable CNF (ECNF)}.
\begin{definition}{Equisatisfiable CNF (ECNF)}{Equisatisfiable CNF (ECNF)}
An \textbf{Equisatisfiable CNF (ECNF)} is a CNF formula that is not logically equivalent to the original formula, but is 
\textbf{Equisatisfiable} -- it has the same satisfiability status:
\begin{itemize}
    \item The original formula is satisfiable if and only if the ECNF is satisfiable 
    \item The ECNF may contain additional auxiliary variables not present in the original formula
\end{itemize}
\end{definition}
\subsubsection{Tseitin Encoding}
The standard technique for producing ECNF is the \textbf{Tseitin encoding} (also called \textbf{Tseitin transformation}),
named after Grigori Tseitin.
\vspace{1em}
\noindent
\newline
\textbf{Key idea:} Introduce auxiliary variables to represent subformulas, avoiding the need for distribution.
\begin{theorem}{Tseitin Encoding Complexity}{Tseitin Encoding Complexity}
\textit{The Tseitin encoding produces an ECNF formula whose size is \textbf{linear} in the size of the original formula: 
O(n) where n is the size of the input formula}
\end{theorem}
\textbf{Advantages of ECNF}:
\begin{itemize}
    \item \textbf{Linear size}: Avoid exponential blowup 
    \item \textbf{Preserves satisfiability}: The ECNF is satisfiable if and only if the original formula is satisfiable 
    \item \textbf{Efficient encoding}: Can be computed in polynomial time 
    \item \textbf{Practical for planning}: Makes SAT-based planning feasible for large problems
\end{itemize}
\textbf{Trade-off}:
\begin{itemize}
    \item The ECNF introduces additional variables (one for each subformula)
    \item The ECNF is not logically equivalent to the original formula (only equisatisfiable)
    \item However, the satisfying assignments to the original variables correspond to solutions of the original formula
\end{itemize}
\subsection{Tseitin Encoding: Formal Construction}
The Tseitin encoding is based on two key functions that work together to produce the ECNF.
\subsubsection{The Representative function}
\begin{definition}{Representative function}{Representative function}
Let PL denote the language of propositional formulas over a set of propositional variables \( \mathcal{X} \), and 
let \( \mathcal{V} \) be a set of fresh (new) propositional variables. The \textbf{representative function} is defined as:
\[
  \text{Rep } : \text{ PL } \rightarrow \mathcal{V} \cup \left\{ \bot, \top \right\}
\]
\end{definition}
The function Rep assigns to each formula a representative variabe (or constant). It is defined inductively:
\vspace{1em}
\noindent
\newline
\textbf{Base Case} (atomic formulas):
\begin{align*}
    & \text{Rep}(x) & = & x \text{ for all } x \in \mathcal{X} \\
    & \text{Rep}(\bot) & = & \bot \\
    & \text{Rep}(\top) & = & \top \\
\end{align*}
\textbf{Inductive Case} (complex formulas):
\[
  \text{Rep}(F) = v_{F}
\]
where \( F \in \text{ PL} \) is a non-atomic formula (i.e., \( F \) is neither a propositional variable nor \( \bot \) nor \( \top \)),
and \( v_F \in \mathcal{V} \) is a fresh propositional variable uniquely associated with $F$.
\vspace{1em}
\noindent
\newline
\textbf{Intuition}: Each complex subformula is represented by a new auxiliary variable. This avoids the need 
to duplicate subformulas during the transformation.
\subsubsection{The Encoding Function}
\begin{definition}{Encoding Function}{Encoding Function}
The \textbf{Encoding function} is defined as:
\[
  \text{En } : \text{ PL } \rightarrow \text{ PL}
\]
\[
  \text{En}(F) = F \leftrightarrow \text{ Rep}(F)
\]
\end{definition} 
This encoding function creates a biconditional formula stating that formula \( F \) is logically equivalent to its 
representative.
\vspace{1em}
\noindent
\newline
\textbf{Interpretation}: \( \text{En}(F) \) encodes the constraint that the auxiliary variable \( \text{Rep}(F) \) must have the same 
truth value as the formula \( F \) itself.
\subsubsection{Complete ECNF Construction}
\begin{definition}{ECNF via Tseitin Encoding}{ECNF via Tseitin Encoding}
Given an input formula \( F \), let \( S_F \) denote the set of all subformulas of \( F \) (including $F$ itself).
The \textbf{Equisatisfiable CNF} of \( F \) is:
\[
  \text{ECNF}(F) = \text{ Rep}(F) \bigwedge_{G \in S_F} \text{ En}(G)
\]
\end{definition}
\textbf{Components}:
\begin{itemize}
    \item \( \text{Rep}(F) \): The representative variable for the entire formula (must be true)
    \item \( \bigwedge_{G \in S_F} \text{En}(G) \): Conjunction of encoding constraints for all subformulas
\end{itemize}
\textbf{Why "Equisatisfiable"?}
\vspace{1em}
\noindent
\newline
The transformation from \( F \) to \( \text{ECNF}(F) \) does \textbf{not} preserve logical equivalence, but it \textbf{does}
preserve satisfiability:
\begin{theorem}{Equisatisfiability Theorem}{Equisatisfiability Theorem}
\textit{Formula F is satisfiable if and only ig ECNF(F) is satisfiable}
\end{theorem}
\textit{Proof Sketch}. \textbf{Forward direction} ($F$ satifiable $\Rightarrow \text{ECNF}(F)$ satisfiable): If $F$ is satisfiable,
we can extend the satisfying assignment to include the auxiliary variables by setting each \( v_G = \text{value}(G) \) under 
the original assignment. This satisfies all encoding constraints.
\vspace{1em}
\noindent
\newline
\textbf{Backward direction} ($\text{ECNF}(F)$ satisfiable $\Rightarrow F$ satisfiable): If $\text{ECNF}(F)$ is satisfiable, the encoding 
constraints ensure that \( \text{Rep}(F) \) has the same truth value as \( F \). Since \( \text{Rep}(F) \) must be 
true (it appears as a conjunct), \( F \) must also be true under the restriction of the assignment to the original variable \( \square \).
\vspace{1em}
\noindent
\newline
\textbf{Key insight}: The ECNF introduces auxiliary variables but maintains the same satisfiability status as the original formula,
while keeping the formula size linear.
\subsection{Example: Equisatisfiability vs. Logical Equivalence}
To illustrate the difference between equisatisfiability and logical equivalence, consider a concrete example.
\begin{example}{Tseitin Encoding of Disjunction}{Tseitin Encoding of Disjunction}
Consider the formula:
\[
  T = A \vee (B \wedge C)
\]
where \( A, B \), and \( C \) are propositional variables.
\end{example}
Using the standard CNF transformation (distribution), we would obtain:
\[
  F_{\text{CNF}} = (A \vee B) \wedge (A \vee C)
\]
This is \textbf{logically equivalent} to \( F \), but for larger formulas, this approach leads to exponential blowup.
\subsubsection{Tseitin Encoding Approach}
Instead, we apply the Tseitin encoding by introducing an auxiliary variable \( v \) to represent the subformula \( B \wedge C \):
\vspace{1em}
\noindent
\newline 
\textbf{Step 1}: Define the encoding constraint:
\[
  v \leftrightarrow (B \wedge C)
\]
\textbf{Step 2}: Replace the subformula in the original formula:
\[
  A \vee v 
\]
\textbf{Step 3}: Convert the biconditional to CNF. The biconditional \( v \leftrightarrow (B \wedge C) \) expands to:
\[
  (v \rightarrow (B \wedge C)) \wedge ((B \wedge C) \rightarrow v) 
\]
Converting each implication:
\begin{align*}
    & v \rightarrow (B \wedge C) \equiv \lnot v \vee (B \wedge C) \equiv (\lnot v \vee B) \wedge (\lnot v \vee C) \\
    & (B \wedge C) \rightarrow v \equiv \lnot(B \wedge C) \vee v \equiv \lnot B \vee \lnot C \vee v \\
\end{align*}
\textbf{Final ECNF}: The complete set of clauses is:
\[
  S = \left\{(A \vee v), (\lnot v \vee B), (\lnot v \vee C), (\lnot B \vee \lnot C \vee v) \right\}
\]
Numbering the clauses:
\begin{align*}
    & (1) & A \vee v \\
    & (2) & \lnot v \vee B \\
    & (3) & \lnot v \vee C \\
    & (4) & \lnot B \vee \lnot C \vee v \\
\end{align*}
\subsubsection{Demonstrating Equisatisfiability}
Consider the interpreation:
\[
  I = \left\{ A \mapsto \text{ true}, B \mapsto \text{ false}, C \mapsto \text{ false} \right\}
\]
\textbf{Evaluating \( F \) under \( I \)}:
\[
  F = A \vee (B \wedge C) = \text{ true } \vee (\text{false } \wedge \text{ false }) = \text{ true } \vee \text{ false } = \text{ true}
\]
Therefore \( I \models F \) (Interpretation $I$ satisfies $F$).
\vspace{1em}
\noindent
\newline 
\textbf{Evaluating \( S \) under \( I \)}:
\vspace{1em}
\noindent
\newline 
We need to extend \( I \) to assign a value to the auxiliary variable \( v \). Let's first try without extending;
\begin{itemize}
    \item Clause (1): \( A \vee v \) = true \( \vee v \) = true \( \quad \) (satisfied regardless of $v$)
    \item Clause (2): \( \lnot v \vee B = \lnot v \vee \) false = \( \lnot v \) \( \quad \) (satisfied only if $v =$ false)
    \item Clause (3): \( \lnot v \vee C = \lnot v \vee v \vee \) false = \( \lnot v \quad \) (satisfied only if $v =$ false)
    \item Clause (4): \( \lnot B \vee \lnot C \vee v = \) true \( \vee \) true \( \vee v \) = true \( \quad \) (satisfied regardless of $v$) 
\end{itemize}
For \( I \) to satisfy \( S \), we must set \( v = \) false (to satisfy clauses (2) and (3)). With this extension:
\[
  I' = \left\{ A \mapsto \text{ true}, B \mapsto \text{ false}, C \mapsto \text{ false}, v \mapsto \text{ false} \right\}
\]
Now \( I' \models S \) (all clause are satisfied).
\begin{definition}{Equisatisfiability vs Logical Equivalence}{Equisatisfiability vs Logical Equivalence}
\textit{Observation} (Not Logically Equivalent). The original interpretation \( I \) (without the auxiliary variable) satisfies 
\( F \) but does \textbf{not} satisfy all clauses in \( S \) unless extended with the appropriate value for \( v \). 
This demonstrates that \( F \) and \( S \) are \textbf{not logically equivalent}. However, \( F \) is satisfiable 
if and only if \( S \) is satisfiable (by appropriately choosing the value of $v$). Therefore, \( F \) and \( S \) are \textbf{equisatisfiable}.
\end{definition}
\textbf{Conclusion}: The Tseitin encoding produces a formula that is equisatisfiable with the original but not logically equivalent.
This is acceptable for SAT-based planning because we only care about whether a solution exists, not about preserving all logical properties.
\subsection{Detailed Definition of the Encoding Function}
We now provide the complete formal definition of the encoding function En for all cases.
\begin{definition}{Encoding Function - Complete Definition}{Encoding Function - Complete Definition}
The encoding function \( \text{En } : \text{ PL } \rightarrow \text{ PL} \) is defined inductively as follows:
\vspace{1em}
\noindent
\newline 
\textbf{Base cases} (atomic formulas):
\begin{align*}
    &\text{En}(\bot) = \top \\
    &\text{En}(\top) = \top \\
    &\text{En}(x) = \top \quad \text{ for all } x \in X 
\end{align*}
\textbf{Inductive case for negation}:
\begin{align*}
    & \text{En}(\lnot F) & = \text{ let } v = \text{ Rep}(\lnot F) \text{ in } v \leftrightarrow \lnot \text{Rep}(F) \\
    & & = (v \rightarrow \lnot \text{Rep}(F)) \wedge (\lnot \text{Rep}(F) \rightarrow v)\\
    & & = (\lnot v \vee \lnot \text{Rep}(F)) \wedge (\text{Rep}(F) \vee v)\\
\end{align*}
\textbf{Inductive cases for binary connectives}:
\vspace{1em}
\noindent
\newline 
For all binary connectives \( \circ \in \left\{ \wedge, \vee, \rightarrow, \leftrightarrow \right\} \):
\[
  \text{En}(F_1 \circ F_2) = \text{ let } v = \text{ Rep}(F_1 \circ F_2) \text{ in } v \leftrightarrow (\text{Rep}(F_1) \circ \text{ Rep}(F_2))
\]
\end{definition}
\textbf{Base cases} (atomic formulas):
\begin{align*}
    & \text{En}(\bot) & = & \top \\
    & \text{En}(\top) & = & \top \\
    & \text{En}(x) & = & \top \quad \text{for all } x \in \mathcal{X} \\
\end{align*}
\textbf{Inductive case for negation}:
\begin{align*}
    & \text{En}(\lnot F) & = \text{ let } v = \text{ Rep}(\lnot F) \text{ in } v \leftrightarrow \lnot \text{Rep}(F)\\
    & & = (v \rightarrow \lnot Rep(F)) \wedge (\lnot\text{Rep}(F) \rightarrow v)\\
    & & = (\lnot v \vee \lnot \text{Rep}(F)) \wedge (\text{Rep}(F) \vee v)\\
\end{align*}
\textbf{Inductive case for binary connectives}:
\vspace{1em}
\noindent 
\newline
For all binary connectives \( \circ \in \left\{ \wedge, \vee, \rightarrow, \leftrightarrow \right\} \):
\[
    \text{En}(F_1 \circ F_2) = \text{ let } v = \text{ Rep}(F_1 \circ F_2) \text{ in } v \leftrightarrow (\text{Rep}(F_1) \circ \text{ Rep}(F_2))
\]
\subsubsection{Example: Conjunction Encoding}
\begin{example}{Conjunction Encoding}{Conjunction Encoding}
Consider the encoding of \( F_1 \wedge F_2 \):
\begin{align*}
    & \text{En}(F_1 \wedge F_2) & = \text{ let } v = \text{ Rep}(F_1 \wedge F_2) \text{ in } v \leftrightarrow (\text{Rep}(F_1) \wedge \text{ Rep}(F_2)) \\
    & & = (v \rightarrow (\text{Rep}(F_1) \wedge \text{ Rep}(F_2))) \wedge ((\text{Rep}(F_1) \wedge \text{ Rep}(F_2)) \rightarrow v) \\
    & & = (\lnot v \vee (\text{Rep}(F_1) \wedge \text{ Rep}(F_2))) \wedge ((\text{Rep}(F_1) \wedge \text{ Rep}(F_2)) \vee v) \\
\end{align*}
Distributing to obtain CNF:
\[
  (\lnot v \vee \text{ Rep}(F_1)) \wedge (\lnot v \vee \text{ Rep}(F_2)) \wedge (\text{Rep}(F_1) \vee v) \wedge (\text{Rep}(F_2) \vee v)
\]
This produces \textbf{4 clauses} in CNF.
\end{example}
\subsection{Complexity Analysis of Tseitin Encoding}
We now analyze the size of the resulting ECNF formula to prove linear complexity.
\begin{theorem}{Linear Size of ECNF}{Linear Size of ECNF}
\textit{Let F be a propositional formula with size |F| = n. The Tseitin encoding produces an ECNF formula:}
\[
  F' = Rep(F) \wedge \bigwedge_{G \in S_F} En(G)
\]
\textit{where \( S_F \) is the set of all subformulas of F (including F itself). The size of F' is |F'| = O(n)}.
\end{theorem}
\subsubsection{Worst Case Analysis: Biconditional}
The worst case for the encoding function occurs with the biconditional connective. Let us analyze:
\begin{example}{Biconditional Encoding - Complete Derivation}{Biconditional Encoding - Complete Derivation}
We show the complete step-by-step conversion of the biconditional encoding to CNF
\noindent\textbf{Step 1: Start with the encoding}
\[
\En(F_1 \leftrightarrow F_2) = \textbf{let } v = \Rep(F_1 \leftrightarrow F_2) \textbf{ in } v \leftrightarrow (\Rep(F_1) \leftrightarrow \Rep(F_2))
\]

\noindent\textbf{Step 2: Expand the inner biconditional}
\[
= v \leftrightarrow \big[(\Rep(F_1) \to \Rep(F_2)) \land (\Rep(F_2) \to \Rep(F_1))\big]
\]

\noindent\textbf{Step 3: Expand the outer biconditional}
\begin{align*}
= \;& v \to \big[(\Rep(F_1) \to \Rep(F_2)) \land (\Rep(F_2) \to \Rep(F_1))\big] \\
\land\; & \big[(\Rep(F_1) \to \Rep(F_2)) \land (\Rep(F_2) \to \Rep(F_1))\big] \to v
\end{align*}

\noindent\textbf{Step 4: Convert implications to disjunctions}
\begin{align*}
= \;& \neg v \lor \big[(\Rep(F_1) \to \Rep(F_2)) \land (\Rep(F_2) \to \Rep(F_1))\big] \\
\land\; & \neg\big[(\Rep(F_1) \to \Rep(F_2)) \land (\Rep(F_2) \to \Rep(F_1))\big] \lor v
\end{align*}

\noindent\textbf{Step 5: Expand inner implications}
\begin{align*}
= \;& \neg v \lor \big[(\neg\Rep(F_1) \lor \Rep(F_2)) \land (\neg\Rep(F_2) \lor \Rep(F_1))\big] \\
\land\; & \neg\big[(\neg\Rep(F_1) \lor \Rep(F_2)) \land (\neg\Rep(F_2) \lor \Rep(F_1))\big] \lor v
\end{align*}

\noindent\textbf{Step 6: Distribute $\neg v$ over the conjunction (first part)}
\[
= (\neg v \lor \neg\Rep(F_1) \lor \Rep(F_2)) \land (\neg v \lor \neg\Rep(F_2) \lor \Rep(F_1))
\]

\noindent\textbf{Step 7: Apply De Morgan's law to the second part}
\[
\land\; \big[\neg(\neg\Rep(F_1) \lor \Rep(F_2)) \lor \neg(\neg\Rep(F_2) \lor \Rep(F_1)) \lor v\big]
\]

\noindent\textbf{Step 8: Apply De Morgan's law again}
\[
\land\; \big[(\Rep(F_1) \land \neg\Rep(F_2)) \lor (\Rep(F_2) \land \neg\Rep(F_1)) \lor v\big]
\]

\noindent\textbf{Step 9: Distribute disjunction over conjunction}
\begin{align*}
\land\;& (\Rep(F_1) \lor \Rep(F_2) \lor v) \land (\Rep(F_1) \lor \neg\Rep(F_1) \lor v) \\
\land\;& (\neg\Rep(F_2) \lor \Rep(F_2) \lor v) \land (\neg\Rep(F_2) \lor \neg\Rep(F_1) \lor v)
\end{align*}

\noindent\textbf{Step 10: Simplify (remove tautologies)}

\noindent The clauses $(\Rep(F_1) \lor \neg\Rep(F_1) \lor v)$ and $(\neg\Rep(F_2) \lor \Rep(F_2) \lor v)$ are tautologies and can be removed.

\bigskip
\noindent\textbf{Final CNF:} The encoding produces the following clauses:
\begin{align*}
& (\neg v \lor \neg\Rep(F_1) \lor \Rep(F_2)) \\
& (\neg v \lor \neg\Rep(F_2) \lor \Rep(F_1)) \\
& (\Rep(F_1) \lor \Rep(F_2) \lor v) \\
& (\neg\Rep(F_2) \lor \neg\Rep(F_1) \lor v)
\end{align*}
This produces \textbf{4 clauses} after simplification (or up to 8 clauses before removing tautologies).
\vspace{1em}
\noindent
\newline
\textbf{Proof of Linear Complexity}. \textbf{Key observations}:
\begin{enumerate}
    \item The number of subformulas \( |S_F| \) is at most \( n \) (the size of $F$)
    \item Each subformula \( G \in S_F \) contributes at most a constant number of clauses (at most 8 for biconditional)
    \item Therefore, the total number of clauses is at most \( 8n \)    
    \item Each clause has at most 3 literals 
    \item Total size: \( |F'| \leq 8n \times 3 = 24n = O(n) \)    
\end{enumerate}
More precisely, we can bound the size as \( |F'| = O(30n + 2) \), which is still linear in \( n \)     
\vspace{1em}
\noindent
\newline
\textbf{Conclusion}: The Tseitin encoding achieves linear size \( O(n) \) instead of the exponential size \( O(2^n) \) of the standard 
CNF transformation, making it practical for large planning problems.
\end{example}
\part{SAT Decision Procedures}
\section{Decision Procedures for Satisfiability}
We now present a \textbf{decision procedure} for determining the satisfiability of a set clauses. This procedure is fundamental to understanding 
how SAT solvers work and forms the algorithmic foundation for planning via SAT encoding.
\begin{definition}{SAT Decision Problem}{SAT Decision Problem}
Given a set of clauses \( S \) over a set of Boolean variables \( \mathcal{X} \), determine whether tehre exists a model (satisfying assignment)
\( v : \mathcal{X} \rightarrow \left\{ \text{true, false} \right\} \) such that \( v \models S \).
\end{definition}
\subsection{Relation to Constraint Solving}
The SAT problem is closely related to \textbf{Constraint Solving Problems (CSP)} with Boolean variables. It represent 
one of the most fundamental problems in algorithmic theory and was the first problem proven to be \textbf{NP-complete}.
\fbox{%
    $\displaystyle 
    S\ \text{(set of clauses)}
    \;\xrightarrow{\;\text{SAT Decision Procedure}\;}\;
    \left\{
        \begin{array}{l}
            \textbf{YES (Satisfiable)}\\[2pt]
            \exists v : v \models S\\[10pt]
            \textbf{No (Unsatisfiable)}\\[2pt]
            \text{Proof: } S \vdash \square
        \end{array}
    \right.
    $%
}
\section{State Representation and Search Models}
\subsection{Trail-Based State Representation}
The state of the SAT decision procedure is represented using a \textbf{trail} of literals, which maintains 
the current partial assignement of truth values to variables.
\begin{definition}{Trail}{Trail}
A \textbf{trail} \( \varGamma \) is a sequence of literals representing the current partial assignment in the search process.
The empty trail is denoted \( \varepsilon \).
\newline
\noindent
(Search State). A search state is represented as a pair \( \langle S, \varGamma \rangle \) where:
\begin{itemize}
    \item \( S \) is the current set of clauses 
    \item \( \varGamma \) is the current trail (partial assignment)
\end{itemize}
\end{definition}
\subsection{Literal Status in Trails}
Given a trail \( \varGamma \) and a literal \( L \), the status of \( L \) is determined as follows:
\begin{itemize}
    \item \( L \) is \textbf{true} if \( L \in \varGamma \)     
    \item \( L \) is \textbf{false} if \( \lnot L \in \varGamma \) 
    \item \( L \) is \textbf{undefined} if neigher \( L \) nor \( \lnot L \) appears in \( \varGamma \) 
\end{itemize}
\section{Transition Rules}
The SAT decision procedure operates as a \textbf{transition system} with specific rules that transform the search state. There 
are two main types of transitions:
\subsection{Decision Rule}
The \textbf{decision rule} extends the trail by choosing an undefined literal and adding it to the current assignment.
\begin{definition}{Decision Rule}{Decision Rule}
\[
  \langle S \parallel \varGamma \rangle \rightarrow \langle S \parallel \varGamma, L \rangle
\]
where the decision is applicable if:
\begin{itemize}
    \item \( L \notin \varGamma \) (literal not already assigned)
    \item \( \lnot L \notin \varGamma \) (negation not already assigned)
\end{itemize}
\end{definition}
The decision rule represent a \textbf{choice point} in the search process where the algorithm guesses a truth value for an undefined variable.
\subsection{Propagation Rule}
The \textbf{propagation rule} (also called \textbf{unit propagation}) adds literals that are logically implied by the current assignment.
\begin{definition}{Propagation Transition}{Propagation Transition}
\[
  \langle S \parallel \varGamma \rangle \rightarrow \langle S \parallel \varGamma, L \rangle
\]
where the propagation is applicable if there exists a clause \( C = L_1 \vee L_2 \vee \ldots \vee L_{k - 1} \in S \) such that:
\begin{itemize}
    \item \( D = L_1 \vee L_2 \vee \ldots L_{k - 1} \)     
    \item \( \varGamma \models \lnot D \) (all literal in $D$ are false under current assignment)
    \item \( L = L_k \) (the remaining literal becomes implied)
    \item \( L \notin \varGamma \) and \( \lnot L \notin \varGamma \)   
\end{itemize}
\end{definition}
\subsubsection{Justification for implied Literals}
When literal \( L_k \) is propagated with justification clause \( C \), it means:
\begin{itemize}
    \item The clause \( C = L_1 \vee L_2 \vee \ldots \vee L_{k - 1} \vee L_k \) must be satisfied 
    \item All others literals \( L_1, L_2, \ldots, L_{k - 1} \) are currently false under $\varGamma$
    \item Therefore, \( L_k \) must be true to satisfy the clause
\end{itemize}
This is called the \textbf{justification} for the propagation: the clause \( C \) forces \( L_k \) to be true.
\subsubsection{Logical Foundation of Propagation}
The propagation rule is based on the following logical principle:
\vspace{1em}
\noindent
\newline
If we have a clause \( D \vee L \) where \( D = L_1 \vee L_2 \vee \ldots \vee L_{k - 1} \), and we know that \( \lnot D \) (i.e., 
$\lnot L_1 \wedge \lnot L_2 \wedge \ldots \wedge \lnot L_{k - 1}$), then: 
\[
  \lnot D \models L
\]
Because if all the alternative literals are false, the remaining literal must be true to satisfy the disjunction.
\section{Solution States}
The search process terminates when it reaches specific terminal states:
\begin{definition}{Solution States}{Solution States}
\begin{itemize}
    \item \textbf{Success}: The state \( \langle S, \varGamma \rangle \) where all clauses in \( S \) are satisfied by the assignment represented by \( \varGamma \)     
    \item \textbf{Failure}: The state where a conflict is detected (e.g., the empty clause $\square$ is derived)
\end{itemize}
\end{definition}
The decision procedure systematically explores the space of possible assignments using these transition rules until it either finds a satisfying assignment 
or proves that no such assignment exists.
\section{Conflict Analysis and Learning}
When the propagation rule leads to a conflict (derivation of the empty clause $\square$), the SAT solver enters \textbf{conflic analysis mode}
to learn from the conflict and guide future search.
\subsubsection{Conflict-Solving Mode}
\begin{definition}{Conflict State Representation}{Conflict State Representation}
A state in conflict-solving mode is represented as a triple \( \langle S \parallel \varGamma \parallel C \rangle \) where:
\begin{itemize}
    \item \( S \) is the current set of clauses 
    \item \( \varGamma \) is the current trail (partial assignment)
    \item \( C \) is the \textbf{conflict clause} that caused the conflict
\end{itemize}
\end{definition}
The conflict clause \( C \) represents a learned constraint that explains why the current partial assignment leads to a contradiction. This clause 
can be added to the formula to prevent the same conflict in future seach branches.
\subsection{Search Space Complexity}
The SAT decision procedure explores a search space that grows exponentially with the number of variables.
\begin{definition}{Search Space Size}{Search Space Size}
Given a SAT problem with \( n \) Boolean variables, the total number of possible assignments is:
\[
  2^n
\]
where each variable can be assigned either true or false.
\end{definition}
\subsubsection{Search Tree Structure}
The search process forms a binary tree where each level corresponds to a variable assignment. For two variables \( p \) and \( Q \), the search space has the following structure:
\begin{itemize}
    \item \textbf{Level 0}: Start state
    \item \textbf{Level 1}: Variable \( p \) can be assigned 1 or 0
    \item \textbf{Level 2}: For each assignment of \( p \), variable \( Q \) can be assigned 1 or 0
\end{itemize}
This creates four possible complete assignments:
\begin{itemize}
    \item \( p = 1, Q = 1 \)    
    \item \( p = 1, Q = 0 \)    
    \item \( p = 0, Q = 0 \)    
    \item \( p = 0, Q = 1 \)    
\end{itemize}
For \( n \) variables, this creates a complete binary tree of depth \( n \) with \( 2^n \) leaves, representing all possible assignments.
\section{Backtracking Mechanism}
When a conflict is detected or a dead-end is reached, the solver uses \textbf{backtracking} to explore alternative assignments.
\subsection{Backtracking Rule}
\begin{definition}{Backtracking Transition}{Backtracking Transition}
\[
  \langle S \parallel M L^d N \parallel C \rangle \rightarrow \langle S \parallel M \lnot L^d \rangle 
\]
where:
\begin{itemize}
    \item \( M \) is the sequence of literals before the decision 
    \item \( L^d \) is the most recente \textbf{decided literal} (marked with $d$)
    \item \( N \) is the sequence of literals after the decision 
    \item The transition is only applicable if \( N \) contains no decided literals 
\end{itemize}
\end{definition}
Backtracking implements a \textbf{depth-first search strategy} that:
\begin{itemize}
    \item Explores one branch of the search tree completely before trying alternatives 
    \item When a conflict occurs, backtracks to the most recent decision point 
    \item Flips the dicsion (from $L^d$ to $\lnot L^d$) and continues search from there
\end{itemize}
\subsection{Decision vs. Propagated Literals}
The distinction between decided and propagated literals is crucial for backtracking:
\begin{itemize}
    \item \textbf{Decided literals} \( (L^d) \): Choosen by the decision rule, can be flipped during backtracking 
    \item \textbf{Propagated literals}: Forced by unit propagation, cannot be flipped (would immediately cause conflict)
\end{itemize}
The backtracking rule only flips decided literals, ensuring that propagated assignments remain consistent.
\section{Termination and Results}
The search process terminates when it reaches one of two possible outcomes:
\subsection{Success state}
\begin{definition}{Satisfiable Instance}{Satisfiable Instance}
\[
  \langle S \parallel \varGamma \rangle \quad \text{ returns } \quad (\text{"sat"}, \varGamma)
\]
if and only if \( \varGamma \models S \) (i.e., the assignment represented by train $\varGamma$ satisfies all clauses in $S$).
\end{definition} 
When a satisfying assignment is found, the solver returns:
\begin{itemize}
    \item Status: \texttt{"sat"} (satisfiable)
    \item Model: The complete assignment \( \varGamma \) that satisfies the formula
\end{itemize}
\subsection{Failure State}
\begin{definition}{Unsatisfiable Instance}{Unsatisfiable Instance}
\[
  \langle S \parallel \varGamma \parallel C \rangle \quad \text{ returns } \quad \text{"unsat"}
\]
if and only if \( \varGamma \) contains no decision literals (i.e., all literals in the trail were propagated, not decided).
\end{definition}
when the formula is proven unsatisfiable, the solver returns:
\begin{itemize}
    \item Status: \texttt{"unsat"} (unsatisfiable)
    \item The empty clause \( \square \) serves as a certificate of unsatisfiability
\end{itemize}
This occurs when the solver has exhausted all possible assignments without finding a satisfying one, meaning the formula 
is logically inconsistent.
\section{Conflict-Driven Clause Learning (CDCL)}
Modern SAT solvers use \textbf{Conflict -Driven Clause Learning (CDCL)}, an advanced 
technique that combines depth-first search with intelligence learning from conflicts 
to efficently solve large satisfiability problems.
\subsection{CDCL Algorithm Overview}
CDCL implements \textbf{depth-first search strategy} with intelligent backjumping (undo one or  more decision) 
and clause learning:
\begin{definition}{CDCL Components}{CDCL Components}
\textbf{CDCL} combines several key techniques:
\begin{itemize}
    \item \textbf{Unit propagation}: Efficient inference of forced assignments
    \item \textbf{Conflict analysis}: Learning from failures to guide future search 
    \item \textbf{Backjumping}: Intelligent backtracking to erliaer decision levels
    \item \textbf{Clause learning}: Adding learned constraints to prevent repeated conflicts
\end{itemize}
\end{definition}
\subsection{Backjumping: Generalized Backtracking}
Backjumping is a generalization of the basic backtracking (undo one decison) mechanism that allows 
the solver to jump back to erlier decision levels, not just the immediately preceding one.
\begin{definition}{Backjumping vs Backtracking}{Backjumping vs Backtracking}
In \textbf{backjumping}, when a conflict occurs, the solver 
analyzes the conflict cause to determine the earliest decisione level that contributed to conflict.
It then jumps back to that level, rather than just the immediately preceding decision.
\vspace{1em}
\noindent
\newline
\textbf{Advantage}: Backjumping can skip over irrelevant decision levels, significantly reducing 
the search space compared to basic backtracking.
\end{definition}
\subsubsection{Conflict analysis example}
COnsider a conflict clause \( C = L_1 \vee \ldots \vee L_k \) where all literals are currently false. The analysis 
determines:
\begin{itemize}
    \item Why each literal \( L_i \) is false in the current trail \( \varGamma \).
    \item The decision level at which each literal became false
    \item The erliest decision level that contributed to the conflict
\end{itemize}
The solver then backjumps to that earliest relevant decision level, effectively skipping decision 
levels that did not contribute to the conflict.
\section{Binary Resolution}
The theoretical foundation of conflict analysis relies on \textbf{binary resolution}, a foundamental inference rule in propositional logic
\subsection{Resolution Rule}
\begin{definition}{Binary Resolution}{Binary Resolution}
Given two clauses containing complementary literals:
\[
  \frac{C \vee L \quad D \vee \lnot L}{(C \vee D)}
\]
\end{definition}
The \textbf{resolvent} \( C \vee D \) is logically entailed by the two premises. This rule allows deriving new 
clauses that preserve satisfiability.
\subsection{Soundness of Resolution}
Binary resolution is \textbf{sound}, the resolvent is logically entailed by the permise.
\begin{theorem}{Resolution Soundness}{Resolution Soundness}
\( If {C \vee L, D \vee \lnot L} \models R where R is any formula, then {C \vee
L, D \vee \lnot L} \models (C \vee D) \).
\vspace{1em}
\noindent
\newline
\textit{Proof}. Consider any interpretation \( v \):
\begin{itemize}
    \item \( If v \models L \), then since \( v \models D \vee \lnot L \), we must 
        have \( v \models D \). Since \( v \models C \vee L \), we have \( v \models C \vee D \).
    \item \( If v \models \lnot L \) then since \( v \models C \vee L \), we must have \( v \models C \). Since \( v \models D \vee \lnot L \),
        we have \( v \models C \vee D \).
\end{itemize}
In both cases, \( v \models C \vee D \), so the resolvent is logically entailed by the premises.
\end{theorem}
\subsection{Application in Conflict Analysis}
In conflict analysis, resolution is used to explain why conflicts occur and to determine which
decision levels are responsible:
\begin{itemize}
    \item When a literal \( L_k \) in the conflict clause is false due to propagation (not a decision)
    \item The propagation was justified by a clause \( D \vee \lnot L_k \).
    \item By resolving the conflict clause with the justification clause, we can trace back the decision that ultimately caused the conflict.
\end{itemize}
\section{Explanation Mechanism}
The explanation mechenism in CDCL traces the implications that led to a conflict,
allowing the solver to learn more general constraints.
\begin{theorem}{Explanation rule}{Explanation rule}
\[ \langle S \Vert \varGamma \Vert C \rangle \rightarrow \langle S \Vert \varGamma \Vert C \vee D \rangle \]
if \( \lnot L \in \varGamma \) where \( \lnot L \) was propagated using the clause \( D \vee \lnot L \).
\end{theorem}
\subsection{Explanation Process}
The explanation mechanism works as follows:
\begin{enumerate}
    \item Start with a conflict clause \( C = L_1 \vee \ldots \vee L_k \).
    \item For each literal \( L_i \) that is false due to propagation (not a decision):
    \begin{itemize}
        \item Find the justification clause that forced \( \lnot L_i \) into the trail 
        \item Resolve the current conflict clause with the justification clause 
        \item Continue untile all literals in the conflict clause are explained by decision 
    \end{itemize}
    \item The final resolved clause contains only literals that were decided (not propagated)
    \item This clause can be learned and added to the formula to prevent similar conflicts
\end{enumerate}
\part{SAT-Based Techniques}
\section{Modern SAT Solving}
\subsection{CDCL: An evolution of DPLL}
CDCL (\textbf{Conflict-Driven Clause Learning}) represent an evolution of the DPLL decision procedure 
for SAT. While CDCL maintains the same basic decision rule as DPLL, it introduces sophisticated conflict anaysis and learning mechanisms that 
dramatically improve solver efficiency on real-world instances.
\subsubsection{CDCL Transition Rules}
The CDCL procedure employs the following transition rules in search mode:
\begin{definition}{Decision rule}{Decision rule}
\[
  \langle S \Vert \varGamma \rangle \rightarrow \langle S \Vert \varGamma, L^d \rangle
\]
if \( L \notin \varGamma \) and \( \lnot L \notin \varGamma \), where:
\begin{itemize}
    \item \( S \) is the set of clauses 
    \item \( \varGamma \) is the trail (sequence of assigned literals)
    \item \( L^d \) is a decision literal (chosen by heuristic)
\end{itemize}
\end{definition}
\begin{definition}{Propagation rule}{Propagation rule}
\[
  \langle S \Vert \varGamma \rangle \rightarrow \langle S \Vert \varGamma, L_{C \vee L} \rangle 
\]
if \( C \vee L \in S, L \notin \varGamma, \lnot L \notin \varGamma \), and \( \varGamma \models \lnot C \), where:
\begin{itemize}
    \item \( C \) is a disjunction of literals (all false under the current trail)
    \item \( L \) is the unit literal that must be propagated 
    \item The subscript \( C \vee L \) indicates the justification clause for the propagation
\end{itemize}
\end{definition}
\begin{definition}{Success rule}{Success rule}
\[
  \langle S \Vert \varGamma \rangle \rightarrow \text{"sat"}
\]
if \( \varGamma \models S \) (i.e., the trail satisfies all clauses in the set)
\end{definition}
\begin{definition}{Conflict rule}{Conflict rule}
\[
  \langle S \Vert \varGamma \rangle \rightarrow \langle S \Vert \varGamma \Vert C \rangle 
\]  
if \( C \in S \) and \( \varGamma \models \lnot C \) (i.e., the trail conflicts with clause \( C \)), where
\( C \) is the conflict clause that triggers transition to conflict analysis mode.
\end{definition}
This four rules define the behavior of CDCL in \textbf{search mode}, which is identical to the DPLL
procedure. The key difference lies in how CDCL handles conflicts through its conflict analysis 
and learning machanisms.
\subsection{Decision Levels in CDCL}
The trail \( \varGamma \) in CDCL is organized into \textbf{decision levels} that provides
structure ti the search process:
\begin{itemize}
    \item The trail begins at \textbf{decision level 0} (empty trail)
    \item Each decision literal opens a new decision level 
    \item If the maximum decision level in \( \varGamma \) is \( n \), the next decision opens level \( n + 1 \).
    \item All propagations following a decision remain at the same level untile the next decision
\end{itemize}
\subsubsection{Trail as a Stack Structure}
The trail can be conceptualized as a \textbf{stack structure}:
\begin{enumerate}
    \item Initially, the stack is empty (decision level 0)
    \item A decision literal is puhsed onto the stack, opening a new level
    \item Unit propagations following the decision are pushed onto the stack at the same level 
    \item Another decision is maade, opening the next level 
    \item The process continues with alternating decision and propagations
    \item During backtracking/backjumping, literals are popped from the stack until reaching the target decision level
\end{enumerate}
This organization enables efficient backjumping by allowing the solver to quickly identify and 
return to relevant decision levels during conflict analysis.

\subsection{Illustrative example}
\begin{example}{Instance}{Instance}
\( S = \left\{ P, \lnot P \vee Q \vee R, T \vee \lnot Q, \lnot R \vee Q \right\} \)
\end{example}
The inital propagation and decision sequence proceeds as follows·
\[
  \langle S \Vert \epsilon \rangle \xrightarrow{prop} \langle S \Vert P_P \rangle
\]



!!! DA QUI INIZIA LA PARTE DEL GOAT FARINELLI QUINDI RECUPERARE TUTTA LA PARTE PRECEDENTE

\part{Planning Module}
\section{Motivation}
\subsection{Deliberative Action}

\subsubsection{The concept of agent}
\begin{itemize}
    \item Pro-Activity: act to achieve a long term goal 
    \item Rationality: do the right thing 
    \item Adaptation: handle unexpected changes/deviations 
    \item Interaction: other afents, humans.
\end{itemize}
\textbf{What is an agent?}
\vspace{1em}
\noindent
\newline

\subsubsection{Agentic AI}
The same concepts, but now the decision and knowledge are based on LLMs and Transfrormers.
\vspace{1em}
\noindent
\newline
\textbf{What is planning?}
\vspace{1em}
\noindent
\newline
It's based on prediction and search, also use a descriptive models of the actions. 
\begin{itemize}
    \item Predict what the actions will do 
    \item Do not tell how to execute actions 
\end{itemize}
Search over \textbf{Predicted state} and possible orgaization of feasible actions. Different types of actions \( \rightarrow \) different 
predictive models, planning problems and techniques.
3 types of planning are:
\begin{enumerate}
    \item Motion planning 
    \item Perception planning 
    \item Task planning
\end{enumerate}
\subsubsection{Planning and Acting}
\begin{itemize}
    \item Multiple levels of abstraction:
    \begin{itemize}
        \item Actor are organized into physical sybsystems 
        \item deliberation resoning techniques
    \end{itemize}
\end{itemize}
\subsubsection{PLanning, Acting and Learning}
\begin{itemize}
    \item Learnining is a very wide concept
    \begin{itemize}
        \item Train a model vs programming a system 
        \item Learn a reactive function 
        \item Learn a partially specified/potentially wrong model
    \end{itemize}
    \item Key point planning acting and learning should be integrated to relize a robust and effective autonomous system
\end{itemize}
\end{document}